%--------------------------------------------------------------------
\chapter{Varietà e teoria del grado}
\setlength{\parindent}{2pt}

\begin{multicols*}{2}
    \section{Varietà differenziabili e prime definizioni}

    \subsection{Mappe \texorpdfstring{$C^\infty$}{C∞} e diffeomorfismi}

    \begin{definition}[Mappe lisce tra due sottinsiemi] Siano $X \subseteq \RR^k$,
        $Y \subseteq \RR^\ell$ sottinsiemi qualsiasi. Allora una funzione
        $f : X \to Y$ si dice di \textbf{classe $C^\infty$} (o \textit{liscia}) se per ogni
        $x \in X$ esistono un aperto $W_x$ e una funzione $F : W_x \to \RR^\ell$,
        chiamata \textbf{estensione}, di classe $C^\infty$ per cui:
        \[
            \restr{F}{W_x \cap X} = \restr{f}{W_x \cap X}.
        \]
    \end{definition}

    \begin{remark}
        Osserviamo che i sottinsiemi di $X$ della forma $W \cap X$ con $W$ aperto sono esattamente gli
        aperti per la topologia di sottospazio di $X$.
    \end{remark}

    \begin{definition}[Diffeomorfismo]
        Siano $X \subseteq \RR^k$,
        $Y \subseteq \RR^\ell$ sottinsiemi qualsiasi.
        Allora una funzione $f : X \to Y$ si dice \textbf{diffeomorfismo} se
        è un \underline{omeomorfismo}, è liscia e ammette inversa liscia.
    \end{definition}

    \begin{remark}
        Le definizioni di mappa liscia e diffeomorfismo date sono chiaramente compatibili
        con le usuali definizioni date su aperti di $\RR^n$.
    \end{remark}

    \begin{proposition} \label{prop:comp_liscia}
        La composizione di mappe lisce è liscia. La restrizione di una mappa liscia
        è liscia.
    \end{proposition}

    \begin{proof}
        Siano $f : X \to Y$ e $g : Y \to Z$ due mappe lisce con $X \subseteq \RR^k$,
        $Y \subseteq \RR^\ell$, $Z \subseteq \RR^p$. Sia $x \in X$. Allora,
        poiché $f$ è liscia, esistono $W_x \subseteq \RR^k$ aperto e $F : W_x \to \RR^\ell$ liscia tale per cui:
        \[
            \restr{F}{W_x \cap X} = \restr{f}{W_x \cap X}.
        \]
        Analogamente, per $f(x) \in Y$ esistono $U_{f(x)} \subseteq \RR^\ell$ aperto e
        $G : U_{f(x)} \to \RR^p$ liscia tale per cui:
        \[
            \restr{G}{U_{f(x)} \cap Y} = \restr{g}{U_{f(x)} \cap Y}.
        \]
        Pertanto, a meno di restringere $W_x$ per ottenere $F(W_x \cap X) \subseteq U_{f(x)} \cap Y$,
        si ha:
        \[
            \restr{G \circ F}{W_x \cap X} = \restr{g \circ f}{W_x \cap X},
        \]
        dove $g \circ f$ è liscia; questo dimostra che la composizione di mappe
        lisce è liscia. \smallskip

        La restrizione di una mappa è liscia dal momento che è composizione
        di una mappa liscia con una mappa di inclusione, che è liscia in quanto
        si estende all'identità.
    \end{proof}

    \subsection{Varietà differenziabili, varietà chiuse, carte, atlanti, parametrizzazioni locali e funzioni di transizione}

    \begin{definition}[Varietà differenziabile liscia senza bordo]
        Un insieme $M \subseteq \RR^k$ si dice \textbf{varietà (differenziabile liscia senza bordo) di dimensione $m>0$} (o $m$-varietà) se per ogni
        suo punto $x$ esistono un intorno aperto $W_x$ in $\RR^k$ e un diffeomorfismo
        $f_x : W_x \cap M \to U$ verso un aperto $U$ in $\RR^m$. Per $m = 0$, si richiede invece
        che ogni $W_x \cap M$ sia un singoletto. \smallskip

        Le coppie della forma $(f_x, W_x \cap M)$ si dicono \textbf{carte locali}, e formano un \textbf{atlante} della varietà. L'inversa di
        $f_x$ si dice invece \textbf{parametrizzazione locale} di $x$ in $M$.
    \end{definition}

    \begin{definition}[Varietà chiusa]
        Si dice \textbf{varietà chiusa} una varietà (senza bordo) che è compatta.
    \end{definition}

    \begin{remark}
        Le varietà di dimensione zero sono esattamente le unioni di punti isolati.
    \end{remark}

    \begin{remark}
        Le carte locali inducono un ricoprimento aperto di $M$, e quindi, qualora $M$ fosse compatta, si potrebbe
        sempre prendere un atlante finito. \smallskip

        Inoltre, poiché $\RR^k$ è II-numerabile, si può sempre prendere un \underline{atlante numerabile}.
    \end{remark}

    \begin{remark}
        Ogni aperto di $\RR^k$ è una varietà di dimensione $k$. Le superfici sono invece varietà
        di dimensione $2$ immerse in $\RR^3$, le cui parametrizzazioni locali sono indotte dalle parametrizzazioni
        regolari.
    \end{remark}

    \begin{remark}[Gli aperti di varietà sono sottovarietà] \label{rmk:aperti_di_varietà_sono_varietà}
        Se $N$ è un aperto di una $m$-varietà $M$, $N$ eredita da $M$ una
        struttura di $m$-varietà per la quale l'atlante è dato dalle intersezioni
        delle carte locali con $N$ stesso. Infatti $N$ è aperto, e dunque
        l'immagine di una carta locale sarà anch'esso un aperto su $\RR^m$.
    \end{remark}

    \begin{definition}[Funzione di transizione]
        Date due parametrizzazioni locali $f : U \to f(U)$ e $g : V \to g(V)$ con intersezione
        delle immagini \underline{non} vuota, si
        definisce la \textbf{funzione di transizione da $f$ a $g$} come la seguente funzione:
        \[ \boxed{g\inv \circ f : f\inv(g(V)) \to g\inv(f(U)).} \]
    \end{definition}

    \subsection{Prodotto di varietà}

    \begin{proposition}[Prodotto di varietà] \label{prop:prodotto_varietà}
        Siano $M \subseteq \RR^k$ e $N \subseteq \RR^\ell$ varietà di dimensione
        $m$ e $n$. Allora $M \times N \subseteq \RR^{k+\ell}$ è una varietà
        di dimensione $m + n$. \smallskip

        Un atlante per $M \times N$ è dato da:
        \[ \{(f_i \times g_j, (W_i \times Q_j) \cap (M \times N))\}_{i, j}, \]
        dove $\{(f_i, W_i \cap M)\}_i$ è un atlante di $M$
        e $\{(g_i, Q_j \cap N)\}_j$ è un atlante di $N$.
    \end{proposition}

    \begin{proof}
        Segue dal fatto che i prodotti $f_i \times g_j$ sono diffeomorfismi in quanto prodotti
        di diffeomorfismi.
    \end{proof}

    \section{Spazio tangente e differenziale su mappe tra varietà}

    \subsection{Differenziale su aperti di \texorpdfstring{$\RR^n$}{ℝⁿ}}

    Ricordiamo la definizione di differenziale per mappe su aperti di spazi reali:

    \begin{definition}[Differenziale per $f : U \subseteq \RR^k \to \RR^\ell$]
        Sia $U$ un aperto di $\RR^k$ e sia $f : U \to \RR^\ell$ una funzione liscia.
        Allora il \textbf{differenziale $\dif f_x : \RR^k \to \RR^\ell$} nel
        punto $x \in U$ è la funzione tale per cui:
        \[
            \boxed{\dif f_x(h) \defeq \lim_{t \to 0} \frac{f(x+th) - f(x)}{t}.}
        \]
        Equivalentemente $\dif f_x$ è l'unica funzione lineare tale per cui:
        \[
            f(x+h) = f(x) + \dif f_x(h) + o(\norm{h}).
        \]
    \end{definition}

    \begin{remark}
        Il differenziale $\dif f_x$ rispetta alcune proprietà fondamentali:
        \begin{enumerate}[(i.)]
            \item La matrice di $\dif f_x$ è data dallo jacobiano $Jf(x) = (\partial_{x_j} f_i(x))_{i, j}$.
            \item Il differenziale rispetta la \textit{regola della catena} (chain rule):
                  \[ \dif(f \circ g)_x = \dif f_{g(x)} \circ \dif g_x. \]
            \item Data $\id_U$ su un aperto $U \subseteq \RR^k$, allora $\dif (\id_U)_x = \id_{\RR^k}$ per ogni $x \in U$.
            \item Dati $U' \subseteq U \subseteq \RR^k$ con $U'$ aperto, l'inclusione $\iota : U' \to U$ è tale per cui
                  $\dif \iota_x = \id_{\RR^k}$ per ogni $x \in U'$.
            \item Se $L : U \to \RR^\ell$ è lineare, allora $\dif L_x = L$ per ogni $x \in U$.
        \end{enumerate}
    \end{remark}

    \begin{proposition} \label{prop:diffeomorfismo_iso_diff}
        Siano $U \subseteq \RR^k$, $V \subseteq \RR^\ell$ aperti. Sia $f : U \to V$ un diffeomorfismo.
        Allora $k = \ell$ e $\dif f_x$ è un isomorfismo di $\RR^k$ per ogni $x \in U$.
    \end{proposition}

    \begin{proof}
        Sia $g : V \to U$ l'inversa di $f$. Poiché $f$ è un diffeomorfismo, $g$ è liscia. Sia $x \in U$. Allora:
        \begin{enumerate}[(i.)]
            \item $\id_{\RR^k} = \dif (\id_U)_x = \dif (g \circ f)_x = \dif g_{f(x)} \circ \dif f_x$,
            \item $\id_{\RR^\ell} = \dif (\id_V)_{f(x)} = \dif (f \circ g)_{f(x)} = \dif f_x \circ \dif g_{f(x)}$.
        \end{enumerate}
        Da (i.) si deduce che $\dif f_x$ è iniettiva, mentre da (ii.) si deduce che è surgettiva. Dunque
        $\dif f_x$ è un isomorfismo (e quindi vale anche $k = \ell$).
    \end{proof}

    \subsection{Spazio tangente in un punto di una varietà}

    \begin{remark}[Lo spazio tangente è ben definito]
        Sia $M$ una varietà di dimensione $m$. Siano $g : U \to W \cap M$ e
        $h : U' \to W' \cap M$ due parametrizzazioni locali di $x \in M$
        con $g(u) = h(u') = x$. \smallskip

        Supponiamo senza perdita di generalità che $W' = W$ (è sufficiente restringere
        le immagini). La funzione
        $g \circ h\inv$ è un diffeomorfismo in quanto composizione di diffeomorfismi
        (vd. Proposizione \ref{prop:comp_liscia}). Allora per la Proposizione \ref{prop:diffeomorfismo_iso_diff}
        $\dif (h\inv \circ g)_u$ è un isomorfismo. \smallskip

        Osserviamo che:
        \[
            \dif g_u = \dif (h \circ (h\inv \circ g))_u = \dif h_{u'} \circ \dif (h\inv \circ g)_u.
        \]
        Dal momento che $\dif (h\inv \circ g)_u$ è in particolare surgettiva, si ha:
        \[
            \dif g_u(\RR^m) = \dif h_{u'}(\RR^m).
        \]
    \end{remark}

    \begin{definition}[Spazio tangente]
        Sia $x$ un punto di una varietà $M$ di dimensione $m$. Presa una
        parametrizzazione locale $g : U \subseteq \RR^m \to W \cap M$ di
        $x$ con $g(u) = x$, si definisce lo \textbf{spazio tangente di $M$ in $x$} come:
        \[
            \boxed{T_x M \defeq \dif g_u(\RR^m).}
        \]
    \end{definition}

    \begin{proposition}
        Sia $x$ un punto di una varietà $M$ di dimensione $m$. Allora:
        \[
            \boxed{\dim T_x M = m.}
        \]
    \end{proposition}

    \begin{proof}
        Si prenda una parametrizzazione locale
        $g : U \subseteq \RR^m \to W \cap M$ di
        $x$ con $g(u) = x$. È sufficiente dimostrare
        che $\dif g_u$ è una mappa iniettiva. \smallskip

        La mappa $g$ è indotta dalla carta locale tramite
        un'estensione $F : W \to \RR^m$ con:
        \[
            \restr{F}{W \cap M} = g\inv.
        \]
        Osserviamo allora che:
        \[
            \id_{\RR^m} = \dif (F \circ g)_u = \dif F_{x} \circ \dif g_u,
        \]
        da cui si deduce che $\dif g_u$ ammette un'inversa sinistra, ed è
        dunque iniettiva.
    \end{proof}

    \begin{remark}[Spazio tangente in un prodotto di varietà] \label{rmk:tangente_prodotto}
        Siano $M$ e $N$ due varietà di dimensione $m$ e $n$.
        Se $f_i \times g_j$ è una carta locale di $M \times N$,
        come ottenuto nella Proposizione \ref{prop:prodotto_varietà}, allora:
        \begin{align*}
            T_{(m, n)}(M \times N) & = d(f_i^{-1} \times g_j^{-1})_{(m, n)}(\mathbb{R}^{m+n})         \\
                                   & \cong d(f_i^{-1})(\mathbb{R}^m) \times d(g_j^{-1})(\mathbb{R}^n) \\
                                   & = T_m M \times T_n N
        \end{align*}
        Quindi vale il seguente isomorfismo canonico, ottenuto proiettando sulle componenti:
        \[
            \boxed{T_{(m, n)} (M \times N) \cong T_m M \times T_n N.}
        \]
    \end{remark}

    \subsection{Differenziale per mappe lisce tra varietà}

    \begin{remark}[Il differenziale per mappe lisce è ben definito]
        Siano $M \subseteq \RR^k$ una varietà di dimensione $m$, $N \subseteq \RR^\ell$
        un'altra varietà, e sia
        $f : M \to N$ una mappa liscia tra varietà. \smallskip

        Sia $F : W \to \RR^\ell$ un'estensione di $f$ per un intorno aperto di
        $x \in M$. Siano $g : U \subseteq \RR^m \to W \cap M$ una parametrizzazione
        locale di $x$ e $h : V \to N$ una parametrizzazione locale di $f(x)$ con
        $g(u) = x$ e $h(v) = f(x)$.

        \[\begin{tikzcd}
                M && W && N \\
                \\
                U &&&& V
                \arrow["\iota", from=1-1, to=1-3]
                \arrow["F", from=1-3, to=1-5]
                \arrow["g", from=3-1, to=1-1]
                \arrow["{h\inv \circ F \circ g}"', dashed, from=3-1, to=3-5]
                \arrow["h"', from=3-5, to=1-5]
            \end{tikzcd}\]

        Dal diagramma commutativo si deduce che:
        \[
            \dif F_x \circ \dif g_u = \dif h_v \circ \dif (h\inv \circ F \circ g)_u.
        \]
        Pertanto $\dif F_x(T_x M)$ \underline{non} dipende dalla scelta dell'estensione $F$ e vale:
        \[ \dif F_x(T_x M) \subseteq T_{f(x)} N. \]
    \end{remark}

    \begin{definition}[Differenziale su mappe tra varietà]
        Sia $f : M \to N$ una mappa tra varietà. Se $F$ è un'estensione di $f$
        in $x$, si definisce il \textbf{differenziale di $f$ in $x$}
        $\dif f_x : T_x M \to T_{f(x)} N$ come segue:
        \[
            \boxed{\dif f_x \defeq \restr{\dif F_x}{T_x M}.}
        \]
    \end{definition}

    \begin{remark}
        Le proprietà del differenziale su aperti di $\RR^n$ si trasferiscono
        facilmente al differenziale su mappe tra varietà:
        \begin{enumerate}[(i.)]
            \item Il differenziale rispetta la \textit{regola della catena} (chain rule):
                  \[ \dif(f \circ g)_x = \dif f_{g(x)} \circ \dif g_x. \]
            \item Data $\id_M$ per una varietà $M$, allora $\dif (\id_M)_x = \id_{T_x M}$ per ogni $x \in M$.
            \item Dati $M'$ e $M$ sono varietà con $M' \subseteq M$, l'inclusione $\iota : M' \to M$ è liscia,
                  $\dif \iota_x : T_x M' \to T_x M$ è iniettiva e $T_x M'$ è un sottospazio di $T_x M$.
            \item Se $f : M \to N$ è un diffeomorfismo, allora $\dif f_x$ è un isomorfismo per ogni $x \in M$.
        \end{enumerate}
    \end{remark}

    \begin{proposition}[Differenziale per prodotti di varietà] \label{prop:diff_prodotto}
        Sia $f : M \to N \times O$ una mappa liscia, dove $M$, $N$ e $O$ sono varietà.
        Se $f(x) = (g(x), p(x))$, allora $g$ e $p$ sono lisce e vale:
        \[
            \boxed{\dif f_x(h) = (\dif g_x(h), \dif p_x(h)).}
        \]
    \end{proposition}

    \section{Valori regolari e critici}

    \subsection{Prime definizioni}

    \begin{definition}[Punti regolari o critici]
        Sia $f : M \to N$ una mappa liscia tra varietà, con
        $\dim M = m$ e $\dim N = n$. \smallskip

        Sia $x \in M$. Si dice che $x$ è un \textbf{punto critico}
        se $\rk(\dif f_x) < n$, e altrimenti si dice che è un
        \textbf{punto regolare}. \smallskip

        Indichiamo con $\crit(f)$ l'insieme dei punti critici di
        $f$.
    \end{definition}

    \begin{definition}[Valori regolari o critici]
        Sia $f : M \to N$ una mappa liscia tra varietà, con
        $\dim M = m$ e $\dim N = n$. \smallskip

        Sia $y \in N$. Si dice che $y$ è un \textbf{valore critico}
        se è immagine di almeno un punto critico, e altrimenti si dice che
        è un \textbf{valore regolare} (in particolare lo è se
        $f\inv(y) = \emptyset$).
    \end{definition}

    \begin{remark}
        È immediato osservare che l'insieme dei valori critici di
        $f$ è esattamente $f(\crit(f))$.
    \end{remark}

    \begin{remark}[I punti regolari formano un aperto] \label{rmk:punti_regolari_formano_aperto}
        Se $x$ è un punto regolare di una mappa liscia
        $f : M \to N$, esiste sempre un intorno aperto $U$
        di $x$ in $M$ composto di soli punti regolari. \smallskip

        Scelta una parametrizzazione locale $g : U \to g(U)$ di un intorno aperto di $x$,
        si può scegliere infatti una base ``comune'' per ogni $T_y M$ al
        variare di $y$ in $g(U)$, e così si può rappresentare
        $\dif f_y$ matricialmente. \smallskip

        Poiché $x$ è regolare, $\dif f_x$ è surgettiva. Allora
        $\dif f_x$ ammette un minore di taglia massima di
        determinante \underline{non} nullo. Il determinante di questo minore,
        al variare di $y \in g(U)$, varia continuamente; in particolare, per
        il Teorema della permanenza del segno, esiste un intorno di $x$
        in cui continua a essere \underline{non} nullo. \smallskip

        Equivalentemente esiste un intorno aperto di $x$ in cui tutti i punti
        sono regolari.
    \end{remark}

    \subsection{Teorema di invertibilità locale per varietà e lemma della pila di dischi}

    \begin{theorem}[di invertibilità locale per varietà] \label{thm:invertibilità_locale_varietà}
        Siano $M$ e $N$ due varietà di stessa dimensione. Sia
        $f : M \to N$ una mappa liscia. Se $x \in M$ è regolare,
        allora esiste un intorno $A$ di $x$ in $M$ tale per cui
        $\restr{f}{A} : A \to f(A)$ è un diffeomorfismo.
    \end{theorem}

    \begin{proof}
        Sia $g : U \to g(U)$ una parametrizzazione locale di $x$ in $M$
        con $g(u) = x$.
        Sia $h : V \to h(V)$ una parametrizzazione locale di $f(x)$ in $N$
        con $h(v) = f(x)$.
        A meno di restringere il dominio di $g$, possiamo supporre che
        $f(g(U)) \subseteq h(V)$.

        \[\begin{tikzcd}
                M && N \\
                \\
                U && V
                \arrow["f"{description}, from=1-1, to=1-3]
                \arrow["g", from=3-1, to=1-1]
                \arrow["{h\inv \circ f \circ g}"{description}, dashed, from=3-1, to=3-3]
                \arrow["h", from=3-3, to=1-3]
            \end{tikzcd}\]

        Osserviamo che:
        \[ \dif f_x \circ \dif g_u = \dif h_{f(x)} \circ \dif (h\inv \circ f \circ g)_u. \]
        Dal momento che $x$ è regolare, $\dif f_x$ è un isomorfismo. Poiché anche $\dif g_u$
        e $\dif h_{f(x)}$ sono isomorfismi ($g$ e $h$ sono diffeomorfismi, vd. Proposizione
        \ref{prop:diffeomorfismo_iso_diff}), allora anche $\dif (h\inv \circ f \circ g)_u$ è
        un isomorfismo. \smallskip

        Per il Teorema di invertibilità locale sugli aperti di $\RR^n$, allora esiste un
        aperto $A'$ di $U$ su cui $\restr{h\inv \circ f \circ g}{A'}$ è un diffeomorfismo.
        Osserviamo che:
        \[
            \restr{f}{g(A')} = \restr{h}{h\inv(f(g(A')))} \circ \restr{h\inv \circ f \circ g}{A'} \circ \restr{g\inv}{g(A')}.
        \]
        Quindi $\restr{f}{g(A')}$ è un diffeomorfismo in quanto composizione di restrizioni di
        diffeomorfismi (vd. Proposizione \ref{prop:comp_liscia}), e $A \defeq g(A')$ è l'intorno cercato.
    \end{proof}

    \begin{proposition} \label{prop:controimmagine_regolare_finita}
        Sia $M$ una varietà compatta. Sia $f : M \to N$ una mappa liscia tra varietà della
        \underline{stessa dimensione}. Se $y \in N$ è un valore regolare, allora $f\inv(y)$ è un insieme finito.
    \end{proposition}

    \begin{proof}
        Poiché $N$ è T1, $f\inv(y)$ è un chiuso di $M$; pertanto, essendo $M$ compatta, $f\inv(y)$ è
        un compatto. \smallskip

        Mostriamo ora che $f\inv(y)$ è discreto. Sia $x \in f\inv(y)$. Dal momento che $y$ è un valore regolare,
        $x$ è un punto regolare. Quindi esiste per il Teorema \ref{thm:invertibilità_locale_varietà} un
        intorno aperto $A \subseteq M$ di $x$ per cui $\restr{f}{A}$ è un diffeomorfismo. In questo intorno $x$
        è l'unica controimmagine di $y$ mediante $f$, e quindi $f\inv(y)$ è discreto. \smallskip

        Dal momento che $f\inv(y)$ è sia compatto che discreto, $f\inv(y)$ è finito.
    \end{proof}

    \begin{lemma}[della pila dei dischi] \label{lem:pila}
        Siano $M$ e $N$ varietà della \underline{stessa dimensione}.
        Sia $M$ compatta. Sia $f : M \to N$ una mappa liscia con $y \in N$ valore regolare.
        Allora esiste un intorno $V$ di $y$ tale per cui:
        \[
            \abs{f\inv(y')} = \abs{f\inv(y)}, \quad \forall y' \in V.
        \]
    \end{lemma}

    \begin{proof}
        Sappiamo dalla Proposizione \ref{prop:controimmagine_regolare_finita} che
        $f\inv(y)$ è finito. Se $f\inv(y) = \{x_1, \ldots, x_n\}$, riprendendo gli intorni aperti $A_i$
        come definiti nella dimostrazione della Proposizione \ref{prop:controimmagine_regolare_finita},
        allora possiamo definire:
        \[
            V \defeq \bigcap_{i = 1}^n f(A_i) \setminus f\left(M \setminus \bigcup_{i = 1}^n A_i\right).
        \]
        Gli $f(A_i)$ sono aperti dal momento che $f$ è un diffeomorfismo. L'insieme $M \setminus \bigcup_{i = 1}^n A_i$
        è un chiuso in un compatto, e quindi è compatto; allora $f\left(M \setminus \bigcup_{i = 1}^n A_i\right)$ è compatto,
        e dunque chiuso essendo $N$ uno spazio T2. Quindi $f\left(M \setminus \bigcup_{i = 1}^n A_i\right)^c$ è aperto.
        Si conclude dunque che $V$ è aperto. \smallskip

        Su $V$, $\abs{f\inv(-)}$ è necessariamente costante: la prima intersezione assicura che esistano
        almeno $\abs{f\inv(y)}$ controimmagini, e la sottrazione insiemistica assicura che non possano esisterne di più.
    \end{proof}

    \subsection{Misura nulla e teoremi di Sard e Brown}

    \begin{definition}[Sottinsiemi di varietà di misura nulla]
        Sia $A$ un sottinsieme di una varietà $M$ di dimensione $m$.
        Si dice che
        $A$ ha \textbf{misura nulla (rispetto a $M$)} se per ogni
        carta locale $(f, W \cap M)$, $f(A \cap W)$ ha misura nulla
        in $\RR^m$.
    \end{definition}

    \begin{theorem}[di Sard, per le varietà] \label{thm:sard}
        Sia $f : M \to N$ una
        mappa liscia tra due varietà. Allora l'insieme dei valori
        critici $f(\crit(f))$ ha misura nulla in $N$.
    \end{theorem}

    \begin{proof}
        Sia $\{(f_i, W_i \cap M)\}_{i \geq 1}$ un atlante numerabile di $M$ e
        sia $(g, Z \cap N)$ una carta locale di $N$. Poniamo:
        \[
            h_i \defeq g \circ \restr{f}{W_i \cap M} \circ f_i.
        \]
        A meno di restringere o ignorare $W_i$, possiamo
        supporre $f(W_i \cap M) \subseteq Z \cap N$.
        \[\begin{tikzcd}
                {W_i \cap M} && {Z \cap N} \\
                \\
                U && V
                \arrow["{\restr{f}{W_i \cap M}}", from=1-1, to=1-3]
                \arrow["{f_i}"', from=1-1, to=3-1]
                \arrow["g"', from=1-3, to=3-3]
                \arrow["{h_i = g \circ \restr{f}{W_i \cap M} \circ f_i}"', dashed, from=3-1, to=3-3]
            \end{tikzcd}\]
        Osserviamo che:
        \[
            \dif (h_i)_u = \dif g_{f(f_i\inv(u))} \circ \dif f_{f_i\inv(u)} \circ \dif (f_i\inv)_u
        \]
        Allora, poiché $\dif g_{f(f_i\inv(u))}$ e $\dif (f_i\inv)_u$ sono isomorfismi
        ($f_i$ e $g$ sono diffeomorfismi, vd. Proposizione \ref{prop:diffeomorfismo_iso_diff}),
        i valori critici di $f$ sono in corrispondenza con quelli degli $h_i$ tramite $g$. \medskip

        Per il Teorema di Sard sugli aperti di $\RR^n$, $g(f(\crit(f) \cap W_i) \cap Z)$ ha allora
        misura zero. Allora, $g(f(\crit(f)) \cap Z)$, che è un unione numerabile di insiemi di misura
        nulla, ha misura nulla. Quindi $f(\crit(f))$ ha misura nulla per definizione.
    \end{proof}

    \begin{corollary}[di Brown] \label{cor:brown} Sia $f : M \to N$ una
        mappa liscia tra due varietà. Allora l'insieme dei valori
        regolari di $f$ è denso in $N$.
    \end{corollary}

    \begin{proof}
        È sufficiente verificare che in un intorno di un valore critico $y \in N$ ci
        sia almeno un valore regolare. Se così \underline{non} fosse, tramite una carta
        locale si troverebbe la chiusura di un rettangolo di soli valori critici, il
        cui volume è non nullo. Allora $f(\crit(f))$ non avrebbe misura nulla, che
        è assurdo per il Teorema \ref{thm:sard}.
    \end{proof}

    \subsection{Varietà a partire da valori regolari}

    \begin{theorem} \label{thm:varietà_da_valore_regolare}
        Sia $f : M \to N$ una mappa liscia tra varietà con $\dim M = m \geq n = \dim N$. Se
        $y \in N$ è regolare, allora $f\inv(y)$ è una varietà di dimensione $m - n$
        \textnormal{(codimensione $n$)}.
    \end{theorem}

    \begin{proof}
        Sia $k$ tale per cui $M \subseteq \RR^k$, e sia $x \in f\inv(y)$.
        Allora $x$ è per ipotesi un punto regolare, e quindi
        $\dif f_x : T_x M \to T_y N$ è una mappa surgettiva, con:
        \[
            \dim \ker(\dif f_x) = m - n.
        \]
        Sia $L : \RR^k \to \RR^{m-n}$ una mappa lineare, dove $T_x M \subseteq \RR^k$
        e $\restr{L}{T_x M}$ è isomorfismo. \smallskip

        Consideriamo la mappa $F : M \subseteq \RR^k \to N \times \RR^{m-n}$ tale per cui:
        \[
            F(m) = (f(m), L(m)).
        \]
        Allora, per la Proposizione \ref{prop:diff_prodotto}, vale:
        \[
            \dif F_x(v) = (\dif f_x(v), \dif L_x(v)) = (\dif f_x(v), L(v)),
        \]
        dove si è usato che $L$ è una mappa lineare. $\dif F_x(v)$ si annulla solo
        per $v = 0$, essendo $\restr{L}{T_x M}$ un isomorfismo; quindi
        $\dif F_x(v)$ è invertibile, e $x$ è regolare per $F$. \smallskip

        Osserviamo che $F$ è una mappa tra varietà della stessa dimensione (vd. Proposizione
        \ref{prop:prodotto_varietà}), e quindi, per la Proposizione
        \ref{thm:invertibilità_locale_varietà}, esiste un intorno $U$ di $x$ in $M$ tale per cui
        $\restr{F}{U} : U \to V \defeq F(U)$ è un diffeomorfismo. \smallskip

        La restrizione $\restr{F}{U}$ mappa $U \cap f\inv(y)$ su un aperto di $V \cap (\{y\} \times \RR^{m-n})$,
        che è diffeomorfo a un aperto di $\RR^{m-n}$ ($\{y\} \times \RR^{m-n} \cong \RR^{m-n}$). In particolare
        induce una carta locale per $x$, e quindi $f\inv(y)$ è una varietà di dimensione $m-n$.
    \end{proof}

    \begin{proposition} \label{prop:tangente_valore_regolare}
        Sia $f : M \to N$ una mappa liscia tra varietà con $\dim M = m \geq n = \dim N$. Se
        $y \in N$ è regolare, posto $P \defeq f\inv(y)$, si ha:
        \[
            \boxed{T_x P = \ker \dif f_x, \quad \forall x \in P.}
        \]
        Inoltre $\restr{\dif f_x}{(T_x P)^\perp}$ è un isomorfismo.
    \end{proposition}

    \begin{proof}
        Poiché $P = f\inv(y)$, l'inclusione $\iota : P \to M$ è tale
        per cui $f \circ \iota$ è costante. Tuttavia una mappa costante
        ha differenziale nullo, e dunque:
        \[
            \dif f_x \circ \underbrace{\iota_{T_x P}}_{\mathclap{= \, \dif \iota^P_x}} = 0 \implies T_x P \subseteq \ker(\dif f_x).
        \]
        Poiché $\dim T_x P = m - n = \ker(\dif f_x)$, si ottiene
        l'uguaglianza $T_x P = \ker(\dif f_x)$. \smallskip

        Osserviamo che $\dim (T_x P)^\perp = n$ e che $(T_x P)^\perp \cap T_x P = \{0\}$.
        Allora $\restr{\dif f_x}{(T_x P)^\perp}$ è iniettiva, e per uguaglianza dimensionale
        tra $(T_x P)^\perp$ e $T_y N$ si conclude che è un isomorfismo.
    \end{proof}

    \begin{corollary} \label{cor:sn_è_varietà}
        $S^n \subseteq \RR^{n+1}$ è una varietà di dimensione $n$ e $T_x S^n = x^\perp \subseteq \RR^{n+1}$.
    \end{corollary}

    \begin{proof}
        Presa $f : \RR^{n+1} \to \RR$ tale per cui:
        \[
            f(x) \defeq x \cdot x = x_1^2 + \ldots + x_{n+1}^2,
        \]
        si ha $S^n = f\inv(1)$. Osserviamo che:
        \[
            Jf_x = 2x^\top.
        \]
        Dunque l'unico valore critico di $f$ è $0$. Pertanto
        $S^n = f\inv(1)$ è una varietà di dimensione $(n+1)-1 = n$. \smallskip

        Poiché $\dif f_x(h) = Jf(x) \cdot h = 2x^\top h$, per la
        Proposizione \ref{prop:tangente_valore_regolare} vale
        anche $T_x S^n = \ker \dif f_x = x^\perp$.
    \end{proof}

    \begin{corollary}
        $O(n) \subseteq M(n)$ è una varietà di dimensione $\frac{n(n-1)}{2}$.
    \end{corollary}

    \begin{proof}
        Presa $f : M(n) \cong \RR^{n^2} \to S(n) \cong \RR^{\frac{n(n+1)}{2}}$ tale per cui:
        \[ f(A) = AA^\top, \]
        si ha $O(n) = f\inv(I)$. Osserviamo che:
        \[ \dif f_A : M(n) \to S(n), \quad \dif f_A(B) = AB^\top + B A^\top. \]
        Mostriamo che $\dif f_A$ è surgettiva per $A \in f\inv(I) = O(n)$. \smallskip

        Sia $C \in S(n)$ simmetrica. Allora $C$ è uguale
        alla sua parte simmetrica:
        \[
            C = \frac{1}{2} C + \frac{1}{2} C^\top,
        \]
        e quindi, ponendo $\frac{1}{2} C = AB^\top$, si ottiene la seguente
        soluzione a $AB^\top + BA^\top = C$:
        \[
            B = \frac{1}{2} C^TA.
        \]
        Dunque $I$ è un valore regolare, e $O(n)$ è una varietà
        di dimensione $n^2 - \frac{n(n+1)}{2} = \frac{n(n-1)}{2}$.
    \end{proof}

    \section{Varietà con bordo}

    \subsection{Semispazio superiore e varietà con bordo}

    \begin{definition}[Semispazio superiore]
        Si definisce il \textbf{semispazio superiore} $H^n$ in $\RR^n$ come:
        \[
            \boxed{H^n \defeq \{x \in \RR^n \mid x_n \geq 0\}.}
        \]
    \end{definition}

    \begin{remark} \label{rmk:hn_diffeo_rn-1}
        Osserviamo che in modo naturale vale il seguente diffeomorfismo:
        \[
            \boxed{\partial H^n = \{ x \in \RR^n \mid x_n = 0 \} \cong \RR^{n-1}.}
        \]
    \end{remark}

    \begin{definition}[$m$-varietà con bordo]
        Si dice che $M \subseteq \RR^k$ è una \textbf{$m$-varietà con bordo} se
        ogni punto di $M$ ammette un intorno diffeomorfo ad un aperto del semispazio
        superiore $H^n$. Gli intorni e i diffeomorfismi citati formano le
        \textbf{carte locali} della varietà, e le inverse di tali diffeomorfismi
        sono dette \textbf{parametrizzazioni locali}. Analogamente si definiscono
        le \textbf{funzioni di transizione}. \smallskip

        Si dice \textbf{bordo} della varietà $M$ l'insieme dei punti che è immagine
        di un punto di $\partial H^n$ tramite qualche parametrizzazione locale, e
        si indica con $\partial M$.
    \end{definition}

    \begin{remark}
        La definizione data è coerente con la definizione di varietà senza bordo: una
        varietà senza bordo $M$ è esattamente una varietà con bordo $M$ tale per cui
        $\partial M = \emptyset$. \smallskip

        Utilizzeremo dunque indistintamente le due caratterizzazioni.
    \end{remark}

    \subsection{Proprietà del bordo di una varietà con bordo}

    \begin{lemma}[I punti di bordo sono sempre immagini di elementi di bordo] \label{lem:punti_di_bordo}
        Sia $x$ un punto del bordo $\partial M$ di una $m$-varietà con bordo $M$. Se
        $g$ è una parametrizzazione locale di $x$, allora $x$ è immagine di un
        punto di bordo di $H^n$ tramite $g$. Equivalentemente, $x$ è un punto di
        $\partial M$ se e solo se è immagine di un valore di bordo per
        ogni sua parametrizzazione locale.
    \end{lemma}

    \begin{proof}
        Sia $g : U \subseteq H^m \to g(U)$ una parametrizzazione locale di $x$ con
        $g(u) = x$. Poiché $x$ è un punto di $\partial M$,
        allora esiste una parametrizzazione locale $f : V \subseteq H^m \to f(V)$
        di $x$ tale per cui esiste $v \in \partial V = \partial H^m \cap V$ con
        $f(v) = x$. \smallskip

        Se $f = g$, la tesi è dimostrata. Se $f \neq g$ e per assurdo $u \notin \partial U$,
        allora la funzione di transizione $g\inv \circ f$ si restringerebbe a
        un diffeomorfismo tra un aperto di $\RR^m$ diffeomorfo a $\RR^m$ e
        un aperto di $H^m$ diffeomorfo a $H^m$. Tuttavia $\RR^m$ e
        $H^m$ non sono diffeomorfi, \Lightning. Dunque $u \in \partial U$.
    \end{proof}

    \begin{corollary}[Il bordo si trasporta naturalmente tramite parametrizzazione locale] \label{cor:bordo_param_locale}
        Sia $g : U \to g(U)$ una parametrizzazione locale di una $m$-varietà con bordo $M$. Allora:
        \[
            \boxed{g(\partial U) = g(U) \cap \partial M.}
        \]
    \end{corollary}

    \begin{proof}
        L'inclusione $g(\partial U) \subseteq g(U) \cap \partial M$ è ovvia.
        L'inclusione opposta invece è data dal Lemma \ref{lem:punti_di_bordo}.
    \end{proof}

    \begin{proposition} \label{prop:bordo_è_varietà}
        Sia $M$ una $m$-varietà con bordo. Allora $\partial M$ è
        una varietà senza bordo di dimensione $m-1$.
    \end{proposition}

    \begin{proof}
        Sia $x \in \partial M$. Allora esiste una parametrizzazione locale
        $g : U \subseteq H^m \to M$ con $g(u) = x$ e $u \in \partial U = U \cap \partial H^m$. \smallskip

        La restrizione $\restr{g}{\partial U}$ è un diffeomorfismo (vd. Proposizione \ref{prop:comp_liscia}). Ricordiamo che
        $g(\partial U) = g(U) \cap \partial M$ dal Corollario \ref{cor:bordo_param_locale}. Allora, poiché
        $\partial H^m \cong \RR^{m-1}$, possiamo identificare
        $\partial U$ come aperto in $\RR^{m-1}$. Quindi
        $(\restr{g}{\partial U})\inv$ induce una carta locale per $\partial M$,
        e si conclude che $\partial M$ è una varietà di dimensione $m-1$. \smallskip
    \end{proof}

    \subsection{Differenziale e spazio tangente su varietà con bordo}

    \begin{remark}[Il differenziale sul bordo di $H^n$ è ben definito]
        ia $g : U \to \RR^k$ una mappa liscia da un aperto $U \subseteq H^n$.
        Supponiamo $\tilde{g}$ e $\hat{g}$ siano due estensioni di $g$ in
        un intorno aperto di $x \in U \cap \partial H^n$. Supponiamo a meno di restringimento
        che $\tilde{g}$ e $\hat{g}$ condividano lo stesso dominio. \smallskip

        Il differenziale $\dif \tilde{g}_x$ coincide allora con
        $\dif \hat{g}_x$. Sia infatti ${u_i}_{i \geq 0}$ è una successione
        in $H^n \setminus \partial H^n$ con $u_i \to x$. Poiché $\tilde{g}$
        e $\hat{g}$ sono lisce, il differenziale vara con continuità, ovverosia:
        \[
            \dif \tilde{g}_x = \lim_{i \to \infty} \dif \tilde{g}_{u_i} =
            \lim_{i \to \infty} \dif \hat{g}_{u_i} = \dif \hat{g}_x,
        \]
        dove si è usato che sugli $u_i$ i differenziali certamente coincidono,
        potendoci restringere a un aperto in $U$ non intersecante il bordo.
    \end{remark}

    \begin{definition}[Differenziale su $H^n$]
        Sia $g : U \to \RR^k$ una mappa liscia da un aperto $U \subseteq H^n$. \smallskip

        Per $x \in U \setminus \partial H^n$, il differenziale $\dif g_x$ è
        definito come l'usuale differenziale dato dalla restrizione di $g$ a un aperto
        di $\RR^n$. \smallskip

        Per $x \in U \cap \partial H^n$, il differenziale $\dif g_x$ è indotto dal
        differenziale di una qualsiasi estensione $\tilde{g}$ di $g$ in un intorno
        aperto di $x$, ovverosia:
        \[
            \boxed{\dif g_x \defeq \dif \hat{g}_x.}
        \]
    \end{definition}

    \begin{remark}
        Come nel caso di una parametrizzazione locale da un aperto di $\RR^n$,
        anche il differenziale di una parametrizzazione locale di una varietà con
        bordo è iniettiva per motivi analoghi.
    \end{remark}

    \begin{definition}[Spazio tangente per varietà con bordo]
        Sia $M \subseteq \RR^k$ una $m$-varietà con bordo. Sia
        $x$ un punto di $M$. Si definisce allora lo \textbf{spazio tangente
            di $x$ su $M$} come:
        \[
            \boxed{T_x M \defeq \dif g_u(\RR^m),}
        \]
        dove $g$ è una parametrizzazione locale di un intorno di $x$ in $M$
        con $g(u) = x$.
    \end{definition}

    \begin{remark}
        Come per il caso di una varietà senza bordo, si dimostra che il
        differenziale è ben definito. Valgono inoltre ancora le usuali
        proprietà del differenziale, inclusa la regola della composizione
        (\textit{chain rule}).
    \end{remark}

    \begin{remark}
        A partire da queste definizioni, si definiscono in modo analogo i
        concetti di punto regolare/critico e di valore regolare/critico. \smallskip

        Si generalizza facilmente in questo senso
        il Teorema di Sard (Teorema \ref{thm:sard}), così come quello di
        Brown (Corollario \ref{cor:brown}).
    \end{remark}

    \begin{remark}[$T_x \partial M$ è un iperpiano di $T_x M$]
        Sia $M$ una $m$-varietà con bordo. Grazie alla Proposizione \ref{prop:bordo_è_varietà}
        sappiamo che $\partial M$ è una $(m-1)$-varietà. \smallskip

        Consideriamo l'inclusione $\iota : \partial M \to M$. Chiaramente $\iota$ è una
        mappa liscia tra varietà con differenziale l'inclusione $T_x \partial M \hookrightarrow T_x M$.
        In particolare vale:
        \[
            \boxed{T_x \partial M \subseteq T_x M}
        \]
        per ogni punto $x \in \partial M$.
    \end{remark}

    \subsection{Varietà con bordo da valori regolari}

    \begin{theorem} \label{thm:varietà_bordata_da_valore_regolare_su_R}
        Sia $f : M \to \RR$ una mappa liscia tra varietà, dove $M$ è una $m$-varietà senza bordo. Se
        $0$ è un valore regolare per $f$, allora $\{ f \geq 0 \}$ è una
        $m$-varietà con bordo $f\inv(0) = \{f = 0\}$.
    \end{theorem}

    \begin{proof}
        L'insieme $\{f > 0\}$ è un aperto di $M$ dal momento che $f$ è continua, e quindi
        eredita la struttura di varietà di $M$ (vd. Osservazione \ref{rmk:aperti_di_varietà_sono_varietà}).
        Dunque, conosciamo già le carte locali di un punto $x \in \{f > 0\}$. \smallskip

        Sia $x$ un punto di $\{f = 0\} = f\inv(0)$. Poiché $0$ è un valore regolare, $\dif f_x$ è
        surgettiva, e quindi $\dim \ker \dif f_x = m - 1$. Supponiamo che $k$ sia tale per cui $M \subseteq \RR^k$.
        Allora possiamo costruire un'applicazione lineare $L : \RR^k \to \RR^{m-1}$ tale per cui
        $\restr{L}{\ker \dif f_x}$ è un isomorfismo. \smallskip

        Consideriamo la mappa $F : M \to \RR^{m-1} \times \RR$ tale per cui:
        \[
            F(m) = (L(m), f(m)).
        \]
        Allora $F$ è una mappa liscia tra varietà, il cui differenziale in $x$ è un isomorfismo. Dunque, per il
        Teorema \ref{thm:invertibilità_locale_varietà}, esiste un intorno aperto $U$ di $x$ in $M$ per il quale
        $\restr{F}{U} : U \to V \defeq F(U)$ è un diffeomorfismo. \smallskip

        Tramite $\restr{F}{U}$ si induce allora un diffeomorfismo tra l'aperto $(\RR^{m-1} \times \RR_{\geq 0}) \cap V = H^m \cap V$
        di $H^m$ e l'aperto $F\inv(H^m \cap V) = \{ f > 0 \} \cap U$, tramite il quale il punto $x$ viene
        mappato su $\partial H^m \cap V$. Dunque $\{ f \geq 0 \}$ è una $m$-varietà con bordo $f\inv(0)$.
    \end{proof}

    \begin{remark}
        Chiaramente il Teorema \ref{thm:varietà_bordata_da_valore_regolare_su_R} si generalizza
        a qualsiasi insieme della forma $\{ f \operatorname{op} a \}$ con $\operatorname{op} \in \{\leq, \geq\}$
        e $a$ valore regolare di $f$.
    \end{remark}

    \begin{corollary}
        $D^n$ è una varietà $n$-dimensionale con bordo $S^{n-1}$.
    \end{corollary}

    \begin{proof}
        Sia $f : \RR^n \to \RR$ tale per cui:
        \[
            f(x) \defeq x \cdot x = x_1^2 + \ldots + x_{n}^2.
        \]
        Allora, come visto per il Corollario \ref{cor:sn_è_varietà},
        $1$ è valore regolare di $f$, e quindi $D^n = \{ f \leq 1 \}$ è
        una $n$-varietà con bordo $S^{n-1}$ per il Teorema \ref{thm:varietà_bordata_da_valore_regolare_su_R}.
    \end{proof}

    \begin{lemma} \label{lem:varietà_bordata_da_valore_regolare_caso_particolare}
        Sia $f : H^m \to \RR^n$ con $m > n$ una mappa liscia. \smallskip

        Se $y \in \RR^n$ è un valore regolare
        sia per $f$ che per $\restr{f}{\partial H^m}$ con $f\inv(y) \neq \emptyset$, allora $f\inv(y)$ è una $(m-n)$-varietà con bordo
        $\partial f\inv(y) = f\inv(y) \cap \partial H^m$.
    \end{lemma}

    \begin{proof}
        Sia $x \in f\inv(y)$. Supponiamo valga $x \in \Int(H^m) = H^m \setminus \partial H^m$. Possiamo
        restringerci a un intorno aperto $U$ di $x$ in $\RR^m$, per il quale $\restr{f}{U}$ diventa una mappa tra
        varietà senza bordo. Allora $\restr{f}{U}\inv(y)$ è una $(m-n)$ varietà senza bordo per il
        Teorema \ref{thm:varietà_da_valore_regolare}. Quindi $x$ eredita da questa varietà le carte locali su $f\inv(y)$;
        inoltre $x$ \underline{non} potrà appartenere al bordo di $f\inv(y)$, essendo nell'immagine di una parametrizzazione
        di un aperto di $\RR^{m-n}$. \smallskip

        Supponiamo valga adesso $x \in \partial H^m$. Poiché $f$ è liscia, esistono un intorno aperto $W$
        di $x$ in $\RR^m$ e un'estensione liscia $F : W \to \RR^n$ per cui:
        \[
            \restr{F}{W \cap H^m} = \restr{f}{W \cap H^m}.
        \]
        Dal momento che $\dif F_x = \dif f_x$, e $x$ è un punto regolare per $f$, allora $\dif F_x$ è surgettiva,
        e $x$ è punto regolare anche per $F$. Dal momento che i punti regolari formano un aperto (vd. Osservazione
        \ref{rmk:punti_regolari_formano_aperto}), possiamo supporre, a meno di restringere $W$, che \underline{non}
        vi siano punti critici in $W$. Pertanto, $y$ sarà valore regolare per $F$ e $F\inv(y)$ è dunque una $(m-n)$-varietà
        senza bordo per il Teorema \ref{thm:varietà_da_valore_regolare}. \smallskip

        Sia $\pi : \RR^m \to \RR$ la proiezione tale per cui $x \mapsto x_m$. Allora $\pi$ è una mappa liscia. \smallskip

        Mostriamo che $0$ è un valore regolare per $\restr{\pi}{F\inv(y)}$. Osserviamo innanzitutto che:
        \[
            (\restr{\pi}{F\inv(y)})\inv(0) = F\inv(y) \cap \partial H^m \underbrace{=}_{\mathclap{(W \setminus H^m) \cap \partial H^m = \emptyset}} f\inv(y) \cap \partial H^m.
        \]
        Sia $y^* \in f\inv(y) \cap \partial H^m$. Osserviamo che:
        \begin{enumerate}[(i.)]
            \item $T_{y^*} F\inv(y) = \ker \dif F_{y^*} = \ker \dif f_{y^*}$ per la Proposizione \ref{prop:tangente_valore_regolare}, e questo spazio
                  ha dimensione $m-n$.
            \item $T_{y^*} \partial H^m = \ker \dif \pi_{y^*}$ per la Proposizione \ref{prop:tangente_valore_regolare}.
            \item $\ker \dif (\restr{f}{\partial H^m})_{y^*} = \ker \restr{\dif f_{y^*}}{T_{y^*} \partial H^m}$ ha dimensione
                  $m-n-1$ dal momento che $y$ è per ipotesi un valore regolare di $\restr{f}{\partial H^m}$, dove per l'uguaglianza
                  dei due nuclei si è utilizzato che:
                  \[
                      \restr{f}{\partial H^m} = f \circ \iota^{\partial H^m},
                  \]
                  e successivamente la \textit{chain rule}.
            \item Grazie a (i.), (ii.) e (iii.), possiamo scrivere:
                  \[
                      \ker \dif (\restr{f}{\partial H^m})_{y^*} = T_{y^*} F\inv(y) \cap (\ker \dif \pi_{y^*}).
                  \]
        \end{enumerate}

        Osserviamo che $y^*$ è un punto critico per $\restr{\pi}{F\inv(y)}$ se e solo se $\dif (\restr{\pi}{F\inv(y)})_{y^*}$ è nullo,
        dacché $\pi$ è una mappa in una $1$-varietà.
        Questo accade se e solo se $T_{y^*} F\inv(y) \subseteq \ker \dif \pi_{y^*}$. Tuttavia, se vi fosse questa inclusione, si avrebbe per (iv.)
        $\ker \dif (\restr{f}{\partial H^m})_{y^*} = T_{y^*} F\inv(y)$, che è assurdo dal momento che il primo spazio ha dimensione $m-n-1$ per (iii.)
        e il secondo ha dimensione $m-n$ per (i.). Quindi $0$ è regolare per $\restr{\pi}{F\inv(0)}$. \smallskip

        Dal momento che $0$ è regolare per $\restr{\pi}{F\inv(0)}$, per il Teorema \ref{thm:varietà_bordata_da_valore_regolare_su_R}
        $\{\pi \geq 0\} = f\inv(y) \cap H^m$ è una $(m-n)$-varietà con bordo $\{\pi = 0\} = f\inv(y) \cap \partial H^m \ni x$.
        Quindi $x$ eredita da questa varietà le carte locali su $f\inv(y)$; inoltre $x$ appartiene al bordo di $f\inv(y)$,
        essendo immagine di un punto di bordo tramite una parametrizzazione locale.
    \end{proof}

    \begin{theorem} \label{thm:varietà_bordata_da_valore_regolare}
        Sia $f : M \to N$ una mappa liscia tra varietà, dove $M$ è una $m$-varietà con bordo $\partial M$ \underline{non} vuoto,
        $N$ è una $n$-varietà senza bordo e $m > n$. \smallskip

        Se $y \in N$ è un valore regolare sia per $f$ che per $\restr{f}{\partial M}$ con $f\inv(y) \neq \emptyset$,
        allora $f\inv(y)$ è una $(m-n)$-varietà con bordo $\partial f\inv(y) = f\inv(y) \cap \partial M$.
    \end{theorem}

    \begin{proof}
        Sia $g : U \subseteq H^m \to g(U) \subseteq M$ una parametrizzazione locale di $x \in f\inv(y)$
        con $g(u) = x$. Sia $h : V \subseteq \RR^n \to h(V) \subseteq N$ una parametrizzazione
        locale di $y$ con $g(v) = y$. \smallskip

        \[\begin{tikzcd}
                M && N \\
                \\
                U && V
                \arrow["f"{description}, from=1-1, to=1-3]
                \arrow["g", from=3-1, to=1-1]
                \arrow["{h\inv \circ f \circ g}"{description}, from=3-1, to=3-3]
                \arrow["h", from=3-3, to=1-3]
            \end{tikzcd}\]

        A meno di restringere i domini delle due mappe, possiamo considerare $p = h\inv \circ f \circ g$.
        Allora $v = h\inv(y)$ è regolare per la mappa $p$. Per possiamo
        restringerci a una palla aperta come intorno aperto di $u$ in $p\inv(v)$: se questa è diffeomorfa a $\RR^m$, una carta locale per $u$
        in $p\inv(v)$ è già data dal Teorema \ref{thm:varietà_da_valore_regolare}; se invece è diffeomorfa a
        $H^m$, la carta locale è ereditata tramite il Lemma \ref{lem:varietà_bordata_da_valore_regolare_caso_particolare}. \smallskip

        La carta locale trovata si trasferisce tramite $g$ al punto $x$ di $M$, e applicando il Corollario \ref{cor:bordo_param_locale}
        il bordo della varietà si trasferisce coerentemente a sua volta.
    \end{proof}

    \subsection{Classificazione delle \texorpdfstring{$1$}{1}-varietà, lemma di non retrazione sul bordo e teorema del punto fisso di Brouwer}

    \begin{theorem}[Classificazione delle $1$-varietà con bordo] \label{thm:classificazione_dim_1_generale}
        Una $1$-varietà con bordo è diffeomorfa a unioni disgiunte di copie di $S^1$ e
        di intervalli di $\RR$.
    \end{theorem}

    \begin{proof}
        Si veda l'appendice di [Milnor, J. (1965). \textit{Topology from the Differentiable Viewpoint}. University Press of Virginia].
    \end{proof}

    \begin{corollary}[Classificazione delle $1$-varietà compatte con bordo] \label{cor:classificazione_dim_1}
        Una $1$-varietà compatta con bordo è necessariamente un'unione disgiunta e finita
        di copie di $S^1$ e di intervalli chiusi di $\RR$.
    \end{corollary}

    \begin{proof}
        Discende immediatamente dal Teorema \ref{thm:classificazione_dim_1_generale}
        utilizzando l'ipotesi di compattezza.
    \end{proof}

    \begin{corollary} \label{cor:punti_pari_su_1varietà}
        Una $1$-varietà compatta con bordo ha un numero pari di punti sul bordo
    \end{corollary}

    \begin{proof}
        Deriva immediatamente dal Corollario \ref{cor:classificazione_dim_1}, dal momento che
        il segmento chiuso ha due punti sul bordo e $S^1$ non ne ha nessuno.
    \end{proof}

    \begin{lemma}[di non retrazione sul bordo] \label{lem:non_retrazione}
        Sia $M$ una varietà compatta con bordo $\partial M \neq \emptyset$. Allora \underline{non}
        esistono mappe lisce $f$ da $M$ in $\partial M$ che fissano il bordo, ovverosia
        con $\restr{f}{\partial M} = \id_{\partial M}$.
    \end{lemma}

    \begin{proof}
        Supponiamo esista una tale mappa $f$. Allora per il Teorema di Brown (Corollario \ref{cor:brown}, generalizzato alle varietà
        bordate), sappiamo che i valori regolari di $f$ sono densi in $\partial M$, e che in particolare
        esiste almeno un valore regolare. \smallskip

        Sia $y$ un
        tale valore regolare. Allora, poiché $\restr{f}{\partial M} = \id_{\partial M}$,
        $y$ è valore regolare anche per $\restr{f}{\partial M}$. \smallskip

        Pertanto, per il Teorema \ref{thm:varietà_bordata_da_valore_regolare}, $f\inv(y)$ è una
        $1$-varietà con bordo $f\inv(y) \cap \partial M$. Poiché $\{y\}$ è chiuso in $\partial M$
        (è T1), $f\inv(y)$ è un chiuso in un compatto, e dunque è una $1$-varietà compatta. \smallskip

        Osserviamo che il bordo di $f\inv(y)$ si semplifica a $\{y\}$ dacché $\restr{f}{\partial M} = \id_{\partial M}$. Tuttavia questo
        è un assurdo, dal momento che le $1$-varietà compatte con bordo finito hanno un numero pari di elementi sul
        bordo per il Corollario \ref{cor:punti_pari_su_1varietà}. Quindi $f$ \underline{non} può esistere.
    \end{proof}

    \begin{lemma} \label{lem:punto_fisso_cinf}
        Ogni mappa liscia $f : D^n \to D^n$ ammette almeno un punto fisso.
    \end{lemma}

    \begin{proof}
        Supponiamo per assurdo che $f$ \underline{non} ammetta alcun punto fisso.
        Definiamo $u_x$ in modo tale che:
        \[
            u_x \defeq \frac{x - f(x)}{\norm{x - f(x)}}, \quad \forall x \in D^n.
        \]
        Poiché $f$ non ammette punti fissi, $u_x \neq 0$ per ogni $x \in D^n$. \smallskip

        Costruiamo $f : D^n \to S^{n-1}$ liscia tale per cui $f(x) = x + t u_x$ e
        $\restr{f}{S^{n-1}} = \id_{S^{n-1}}$. \smallskip

        Allora deve valere $\norm{f(x)}^2 = 1$, e quindi:
        \begin{equation*}
            \begin{aligned}
                t^2 + 2(x \cdot u_x) t + (\|x\|^2 - 1) = 0 \implies \\
                t_{\pm} = - x \cdot u_x \pm \sqrt{(x \cdot u_x)^2 - \|x\|^2 + 1}.
            \end{aligned}
        \end{equation*}
        Scegliamo $t \defeq t_+$. Affinché $f$ sia liscia, occorre che l'espressione
        $h(x) \defeq (x \cdot u_x)^2 - \|x\|^2 + 1$ dentro la radice quadrata in $t$ sia sempre maggiore di $0$. \smallskip

        Osserviamo che per $x \in D^n$ si ha $h(x) \geq (x \cdot u_x)^2 \geq 0$. Se allora $h(x) > (x \cdot u_x)^2$,
        $h(x)$ è maggiore strettamente di $0$. Se invece $h(x) = (x \cdot u_x)^2$, necessariamente $\norm{x}^2 = 1$,
        e quindi $x \in S^{n-1}$. Ma allora:
        \[
            x \cdot u_x = \frac{1 - x \cdot f(x)}{\norm{x - g(x)}}.
        \]
        Dal momento che $x$ non può essere un punto fisso di $f$, $x \cdot f(x)$ \underline{non} può essere uguale a $1$; quindi
        $x \cdot u_x > 0$ per $x \in S^{n-1}$. Si conclude allora che $h(x) > 0$ per ogni $x \in D^n$, e quindi $f$
        è una funzione liscia. \smallskip

        Si verifica facilmente che $\restr{f}{S^{n-1}} = \id_{S^{n-1}}$. Questo tuttavia è un assurdo
        per il Lemma \ref{lem:non_retrazione}, \Lightning. Quindi $f$ ammette almeno un punto fisso.
    \end{proof}

    \begin{theorem}[del punto fisso di Brouwer]
        Ogni mappa continua $f : D^n \to D^n$ ammette almeno un punto fisso.
    \end{theorem}

    \begin{proof}
        Supponiamo per assurdo che $f$ \underline{non} ammetta alcun punto fisso. Poiché $D^n$
        è compatto, per il Teorema di approssimazione di Weierstrass, per ogni $\eps > 0$
        esiste una funzione polinomiale $P_\eps : \RR^n \to \RR^n$ tale per cui
        $\norm{P_\eps(x) - f(x)} < \eps$ per ogni $x \in D^n$. \smallskip

        Osserviamo che:
        \[
            \norm{P_\eps(x)} \leq \norm{P_\eps(x) - f(x)} + \norm{f(x)} < \eps + 1,
        \]
        quindi $Q_\eps \defeq \frac{1}{\eps + 1} P_\eps$ si restringe su $D^n$ a
        un'endofunzione $\restr{Q_\eps}{D^n} : D^n \to D^n$. \smallskip

        Dal momento che $f$ non ammette alcun punto fisso, $\mu \defeq \min_{x \in D^n} \norm{f(x) - x}$ è
        tale per cui $\mu > 0$. \smallskip

        Osserviamo che:
        \begin{equation*}
            \begin{aligned}
                \|Q_\eps(x) - f(x)\| & = \frac{1}{\eps + 1} (\|P_\eps(x) - f(x)\| + \eps \|f(x)\|) \\[0.1in]
                                     & \leq \frac{2 \eps}{\eps + 1}.
            \end{aligned}
        \end{equation*}
        Scegliendo allora $\eps$ tale per cui $\frac{2 \eps}{\eps + 1} < \mu$, $Q_\eps$ \underline{non} può ammettere
        punti fissi: un punto fisso violerebbe infatti la minimalità di $\mu$ secondo la scorsa disuguaglianza.
        Tuttavia $Q_\eps$ è una funzione liscia, e per il Lemma \ref{lem:punto_fisso_cinf} deve ammettere
        punti fissi, \Lightning. Dunque $f$ ammette almeno un punto fisso.
    \end{proof}

    \section{Teoria del grado modulo \texorpdfstring{$2$}{2}}

    \subsection{Omotopie \texorpdfstring{$C^\infty$}{C∞}}

    \begin{remark}[{$M \times [0,1]$} è una varietà con bordo] \label{rmk:m_01_varietà}
        Sia $M$ una $m$-varietà senza bordo. Allora, per la Proposizione \ref{prop:prodotto_varietà},
        $M \times \RR$ è una $(m+1)$-varietà senza bordo. \smallskip

        Consideriamo la mappa liscia $f : M \times \RR \to \RR$ tale per cui
        $f(x, t) = t(t-1)$. Allora per il Teorema \ref{thm:varietà_bordata_da_valore_regolare_su_R}
        $\{f \leq 0\} = M \times [0, 1]$ è una $(m+1)$-varietà con bordo
        $\{f = 0\} = M \times \{0\} \sqcup M \times \{1\}$.
    \end{remark}

    \begin{definition}[Omotopia $C^\infty$ e funzioni $C^\infty$-omotope]
        Siano $f$ e $g$ due funzioni da una varietà $M$ in una $N$.
        Un'\textbf{omotopia $C^\infty$} da $f$ a $g$ è una funzione liscia
        $H : M \times [0, 1] \to N$ tale per cui:
        \begin{itemize}
            \item $H(-, 0) = f$,
            \item $H(-, 1) = g$.
        \end{itemize}
        Definiamo inoltre:
        \[
            \boxed{H_t \defeq H(-, t).}
        \]

        Due funzioni $f$ e $g$ per le quali esiste un'omotopia
        da $f$ a $g$ si dicono \textbf{$C^\infty$-omotope}.
    \end{definition}

    \begin{remark}
        È immediato verificare che ``essere $C^\infty$-omotope'' è una
        relazione di equivalenza per le funzioni lisce da $M$ a $N$.
    \end{remark}

    \begin{lemma}[di omotopia] \label{lem:omotopia}
        Siano $f$ e $g$ due funzioni $C^\infty$-omotope da una varietà $M$ in una $N$,
        con $M$ compatta e $\dim M = \dim N$. Se
        $y \in N$ è un valore regolare sia per $f$ che per $g$, allora:
        \[
            \boxed{\abs{f\inv(y)} \equiv \abs{g\inv(y)} \pmod{2}.}
        \]
    \end{lemma}

    \begin{proof}
        Sia $H$ una omotopia $C^\infty$ da $f$ a $g$. Allora, poiché $M$ è compatta,
        per il Lemma \ref{lem:pila}:
        \begin{itemize}
            \item esiste un intorno $V_1$ di $y \in N$ su cui $\abs{f\inv(-)}$ è costante;
            \item esiste un intorno $V_2$ di $y \in N$ su cui $\abs{g\inv(-)}$ è costante.
        \end{itemize}
        È sufficiente allora mostrare la tesi per un qualsiasi valore $y' \in V_1 \cap V_2$;
        poiché i valori regolari sono densi per il Teorema di Brown (per le varietà, Corollario \ref{cor:brown}),
        possiamo prendere $y' \in V_1 \cap V_2$ valore regolare di $H$. \smallskip

        Allora $H\inv(y')$ è una varietà di dimensione $1$, il cui bordo ha un numero pari
        di punti per il Corollario \ref{cor:punti_pari_su_1varietà}. \smallskip

        Osserviamo che $\partial (M \times [0, 1]) = M \times \{0\} \sqcup M \times \{1\}$. Allora:
        \[
            \partial H\inv(y') = (f\inv(y') \times \{0\}) \sqcup (g\inv(y') \times \{1\}).
        \]
        Quindi:
        \[
            0 \equiv \abs{\partial H\inv(y')} \equiv \abs{f\inv(y)} + \abs{g\inv(y)} \pmod{2}.
        \]
    \end{proof}

    \subsection{Isotopie e lemma di omogeneità}

    \begin{definition}[Isotopia]
        Una omotopia $C^\infty$ $H : M \times [0, 1] \to N$ si dice
        \textbf{isotopia} se per ogni $t \in [0, 1]$, $H(-, t)$ è un
        diffeomorfismo liscio.
    \end{definition}

    \begin{definition}[Isotopia a supporto compatto]
        Un'isotopia $H : N \times [0, 1] \to N$ si dice \textbf{a supporto compatto} se esiste
        un compatto $K \subseteq N$ tale per cui $\restr{H(-, t)}{N \setminus K} = \id_{N \setminus K}$
        per ogni $t \in [0, 1]$.
    \end{definition}

    \begin{definition}[Punti isotopi di una varietà]
        Due punti $y$, $z \in N$, dove $N$ è una varietà, si dicono
        \textbf{isotopi} se esiste un diffeomorfismo $h : N \to N$
        con $h(y) = z$ e un'isotopia a supporto compatto da $h$ a
        $\id_{N}$.
    \end{definition}

    \begin{remark}
        È immediato osservare che la relazione ``essere isotopi'' sui punti di una varietà è una relazione di equivalenza.
    \end{remark}

    \begin{lemma} \label{lem:classi_equivalenza_aperte}
        Le classi di equivalenza della relazione ``essere isotopi'' sui
        punti di una varietà (senza bordo) $M$ sono aperti della varietà.
    \end{lemma}

    \begin{proof}
        Sia $x$ un punto di $M$. Supponiamo $M$ sia una $(n+1)$-varietà.
        Usando le carte locali, eventualmente riparametrizzate
        per avere come immagine una palla di centro $0$, è sufficiente mostrare
        che per ogni $r$, esiste $0 < r_0 \leq r$ tale per cui
        tutti i punti di $B_{r_0}(0) \subseteq B_r(0) \subseteq \RR^{n+1}$ sono isotopi a $0$. \smallskip

        Possiamo ridurre il problema ulteriormente concentrandoci sui punti
        della forma $z = (a, 0, \ldots, 0)$ con $a < r_0$, dal momento che una rotazione
        permette poi di rendere isotopi tutti gli altri punti di $\partial B_a(0)$. Se
        infatti $H$ è un isotopia a supporto compatto che porta $z$ in $0$, allora:
        \[ H'(-, t) \defeq R_\theta\inv \circ H(-, t) \circ R_\theta \]
        è l'isotopia cercata, dove $R_\theta$ è l'opportuna rotazione scelta. \smallskip

        Siano $\rho : \RR \to [0, 1]$ e $\sigma : \RR^n \to [0, 1]$ due funzioni di test
        (\textit{bump function}) lisce con $\rho(0) = \sigma(0) = 1$ e:
        \[ \abs{x} > 1 \implies \rho(x) = 0, \quad \norm{y} > 1 \implies \sigma(y) = 0. \]
        Osserviamo che:
        \[
            \abs{x} > \eps \implies \rho\left(\frac{x}{\eps}\right) = 0, \quad \norm{y} > \delta \implies \sigma\left(\frac{y}{\delta}\right) = 0.
        \]

        Scelti allora $\eps$ e $\delta$ di modo che $\eps^2 + \delta^2 < r_0^2$, definiamo allora l'omotopia $C^\infty$ $H$ tale per cui:
        \begin{equation*}
            \begin{aligned}
                H : \overbrace{\mathbb{R} \times \mathbb{R}^n}^{\mathclap{\mathbb{R}^{n+1}}} \times [0, 1] & \to \overbrace{\mathbb{R} \times \mathbb{R}^n}^{\mathclap{\mathbb{R}^{n+1}}},                         \\
                H(x, y, t)                                                                                 & = \left(x + a t \rho\left(\frac{x}{\varepsilon}\right) \sigma\left(\frac{y}{\delta}\right), y\right).
            \end{aligned}
        \end{equation*}

        Osserviamo subito che $H_0 = \id_N$, $H_1(0) = z$ e che $H_t$ fissa tutto ciò che è fuori
        da $B_{r_0}(0)$. Mostriamo che per $r_0$ sufficientemente piccolo $H_t$ è un diffeomorfismo. \smallskip

        Lo jacobiano di $H_t$ è il seguente:
        \[
            J H_t = \begin{pmatrix}
                1 + \frac{a t}{\varepsilon} \sigma\left(\frac{y}{\delta}\right) \rho'\left(\frac{x}{\varepsilon}\right) & \vline & *   \\[0.06in]
                \hline
                0                                                                                                       & \vline & I_n
            \end{pmatrix}.
        \]
        Si verifica facilmente che si può scegliere $r_0$, e successivamente $\eps$ e $\delta$ in modo tale che:
        \[
            \abs{\frac{a t}{\varepsilon} \sigma\left(\frac{y}{\delta}\right) \rho'\left(\frac{x}{\varepsilon}\right)} < 1, \quad \forall a, x, y,
        \]
        e conseguentemente $\det(J H_t) > 0$, da cui si deduce che $H_t$ è un diffeomorfismo locale. Analogamente, $H_t$ è
        bigettiva sulle rette $\{y = \text{cost.}\}$ dal momento che la derivata sulle rette è strettamente
        positiva. Dunque $H_t$ è bigettiva e diffeomorfismo locale, e quindi è diffeomorfismo. \smallskip

        Si conclude che $H$ è l'isotopia cercata, e quindi ogni classe di isotopia è aperta.
    \end{proof}

    \begin{lemma}[di omogeneità] \label{lem:omogeneità}
        Sia $N$ una varietà connessa e siano $y$, $z$ due suoi punti. Allora
        $y$ e $z$ sono isotopi.
    \end{lemma}

    \begin{proof}
        Poiché $N$ è connessa, per il Lemma \ref{lem:classi_equivalenza_aperte} esiste
        allora un'unica classe di equivalenza per la relazione ``essere isotopi'', da
        cui segue immediatamente la tesi.
    \end{proof}

    \subsection{Grado modulo \texorpdfstring{$2$}{2} e buona definizione}

    \begin{theorem}
        Siano $M$ e $N$ varietà della stessa dimensione. Sia $M$ chiusa (ovverosia
        anche compatta) e $N$ connessa. Siano $y$ e $z$ due valori regolari di una
        funzione $f : M \to N$ liscia. Allora:
        \[ \abs{f\inv(y)} \equiv \abs{f\inv(z)} \pmod{2}. \]
    \end{theorem}

    \begin{proof}
        Poiché $N$ è connessa, per il Lemma \ref{lem:omogeneità} esiste un diffeomorfismo
        $h : N \to N$ con $h(y) = z$ e un'isotopia a supporto compatto
        $H : N \times [0, 1] \to N$ da $\id_N$ a $h$.
        \[\begin{tikzcd}
                {M \times [0, 1]} && {N \times [0, 1]} && N
                \arrow["{f \times \id_{[0, 1]}}", from=1-1, to=1-3]
                \arrow["H", from=1-3, to=1-5]
            \end{tikzcd}\]
        Poiché $y$ è regolare per $f$ e $h$ è diffeomorfismo, si deduce che
        $z$ è regolare per $h \circ f$. Consideriamo l'omotopia
        $H' = H \circ (f \times \id_{[0,1]})$. Osserviamo che
        $H'_0 = f$ e che $H'_1 = h \circ f$.
        Quindi, per il Lemma \ref{lem:omotopia}, si conclude che:
        \[
            \abs{f\inv(z)} \equiv \abs{(h \circ f)\inv(z)} \equiv \abs{f\inv(y)} \pmod{2}.
        \]
    \end{proof}

    \begin{definition}[Grado modulo $2$ di una funzione liscia]
        Sia $f : M \to N$ una funzione liscia da una varietà compatta $M$
        a una di stessa dimensione e connessa $N$. Allora si definisce
        il \textbf{grado modulo $2$ di $f$} come:
        \[
            \boxed{\deg_2 f \defeq \abs{f\inv(y)} \bmod 2,}
        \]
        dove $y$ è un qualsiasi valore regolare di $f$.
    \end{definition}

    \begin{lemma} \label{lem:valori_regolari_aperto}
        Siano $M$ e $N$ varietà della stessa dimensione. Sia $M$ chiusa (ovverosia
        anche compatta)
        e $N$ connessa. Se $f : M \to N$ è liscia, allora i valori regolari di $f$
        formano un aperto di $N$.
    \end{lemma}

    \begin{proof}
        Dall'Osservazione \ref{rmk:punti_regolari_formano_aperto} sappiamo che
        i punti regolari di $f$ formano un aperto di $M$; dunque i punti critici
        formano un chiuso di $M$. Dacché $M$ è compatta, in particolare i punti
        critici formano un compatto di $M$. Tramite $f$, si deduce che i valori
        critici formano un compatto di $N$, che, essendo $N$ T2, è dunque
        chiuso in $N$. Quindi i valori regolari formano un aperto.
    \end{proof}

    \begin{theorem} \label{thm:fondamentale_grado_2}
        Siano $M$ e $N$ varietà della stessa dimensione. Sia $M$ chiusa (ovverosia
        anche compatta)
        e $N$ connessa. Se $f$ e $g$ sono due mappe lisce $C^\infty$-omotope da $M$ in $N$,
        allora:
        \[
            \boxed{\deg_2 f = \deg_2 g.}
        \]
    \end{theorem}

    \begin{proof}
        Per il Lemma \ref{lem:valori_regolari_aperto}, i valori regolari di $f$ formano un aperto in $N$. Allora, per il Teorema
        di Brown (per varietà, Corollario \ref{cor:brown}), esiste in questo
        aperto anche un valore regolare di $g$. Per il Lemma \ref{lem:omotopia},
        dunque $f$ e $g$ condividono lo stesso grado modulo $2$.
    \end{proof}

    \begin{corollary}
        La mappa costante $c_{x_0} : S^n \to S^n$ \underline{non} è
        $C^\infty$-omotopa a $\id_{S^n}$.
    \end{corollary}

    \begin{proof}
        Poiché $c_{x_0}$ non è surgettiva, $\deg_2 c_{x_0} = 0$ (ogni punto diverso da $x_0$ è
        valore regolare con controimmagine vuota); mentre
        $\deg_2 \id_{S^n} = 1$. Quindi per il Teorema \ref{thm:fondamentale_grado_2}
        le due mappe non possono essere $C^\infty$-omotope.
    \end{proof}

    \section{Varietà orientate}

    \subsection{Orientazione di basi su spazi vettoriali, orientazione canonica di \texorpdfstring{$\RR^n$}{ℝⁿ}}

    \begin{definition}[Stessa orientazione]
        Si dice che due basi (ordinate) $\basis$, $\basis'$ di un $\RR$-spazio vettoriale finito-dimensionale
        \textbf{hanno la stessa orientazione} se la matrice del cambio di base da $\basis$
        a $\basis'$ ha determinante positivo.
    \end{definition}

    \begin{remark}
        È immediato verificare che ``avere la stessa orientazione'' è
        una relazione d'equivalenza sulle basi dello spazio in esame
        avente solo due classi di equivalenza
        per le orientazioni.
    \end{remark}

    \begin{definition}[Orientazione]
        Definiamo \textbf{orientazione} una classe di equivalenza per la relazione
        ``avere la stessa orientazione''. \smallskip

        Data un'orientazione $\Theta$ indichiamo con $-\Theta$ l'unica altra
        classe di equivalenza.
    \end{definition}

    \begin{definition}[Orientazione canonica]
        Si definisce l'\textbf{orientazione canonica} $\Theta_0$ di $\RR^n$
        come la classe di equivalenza indotta dall'orientazione della base
        canonica. \smallskip
    \end{definition}

    \begin{remark}
        Un isomorfismo $L : V \to V'$ induce una bigezione dalle orientazioni
        di $V$ a quelle di $V'$ tramite:
        \[
            [\basis] \mapsto [L(\basis)].
        \]
        Infatti la matrice di cambio di base è invariante per isomorfismo. \smallskip

        Indicheremo tale mappa con il simbolo dell'isomorfismo da cui è indotta.
    \end{remark}

    \begin{definition}[Segno di una base]
        Dato uno spazio vettoriale $V$ orientato con $\Theta$, si definisce
        il \textbf{segno} di una base $\basis$ come:
        \[
            \boxed{\sgn_\Theta(\basis) \defeq \begin{cases}
                    +1 & \text{se } \basis \sim \Theta, \\
                    -1 & \text{altrimenti}.
                \end{cases}}
        \]
    \end{definition}

    \subsection{Orientazione su prodotti di spazi vettoriali}

    \begin{definition}[Orientazione prodotto]
        Siano $V$ e $W$ due $\RR$-spazi vettoriali di dimensione finita.
        Se $\Theta^V$ è un'orientazione di $V$ e $\Theta^W$ lo è di $W$,
        allora si definisce l'\textbf{orientazione prodotto} $\Theta^V \times \Theta^W$ su
        $V \times W$ come l'orientazione indotta dalla giustapposizione di $\basis_V$ e $\basis_W$:
        \[ \boxed{\basis_V \sqcup \basis_W \defeq \basis_V \times \{0_W\} \cup \{0_V\} \times \basis_W,} \]
        dove $\basis_V$ e $\basis_W$ sono basi di $V$ e $W$ con
        $[\basis_V] = \Theta^V$ e $[\basis_W] = \Theta^W$.
    \end{definition}

    \begin{remark}[Regola dei segni per l'orientazione prodotto] \label{rmk:regola_segni}
        Siano $V$ e $W$ $\RR$-spazi orientati con $\Theta^V$ e $\Theta^W$.
        Siano $\basis_V$ e $\basis_W$ basi di $V$ e $W$. Sia $M_V$
        la matrice di cambio di base da una base positiva di $V$ a $\basis_V$. Sia
        $M_W$ l'analogo per $W$. \smallskip

        La matrice di cambio di base dalla giustapposizione delle basi positive alla giustapposizione di $\basis_V$
        e $\basis_W$ è esattamente:
        \[
            M = \begin{pmatrix}
                M_V & \vline & 0   \\
                \hline
                0   & \vline & M_W
            \end{pmatrix}.
        \]
        Quindi vale la seguente \textit{regola dei segni}:
        \[
            \boxed{\sgn_{\Theta^V \times \Theta^W}(\basis_V \sqcup \basis_W) = \sgn_{\Theta^V}(\basis_V) \sgn_{\Theta^W}(\basis_W).}
        \]
    \end{remark}

    \subsection{Orientazione su varietà e prime proprietà}

    \begin{definition}[$m$-varietà orientata, $m > 1$ o $\partial M = \emptyset$]
        Una \textbf{varietà orientata di dimensione $m$} (con $\underline{m>1}$ o $\underline{\partial M = \emptyset}$)
        è una coppia $(M, \Theta)$, dove $M$ è una $m$-varietà, eventualmente con bordo, e $\Theta = \{\Theta_x\}_{x \in M}$
        è una famiglia di orientazioni degli spazi tangenti dei punti di $M$ tale per cui:

        \begin{quote}
            Per ogni $x \in M$ esiste una parametrizzazione locale $g : U \subseteq \RR^m \to g(U) \subseteq M$ con
            $\dif g_u (\Theta_0) = \Theta_{g(U)}$ per ogni $u \in U$ (\textbf{condizione di compatibilità di $g$ con $\Theta$}).
        \end{quote}

        Una varietà $M$ per cui esiste una famiglia di orientazioni tali per cui $(M, \Theta)$ è orientata si
        dice \textbf{orientabile}.
    \end{definition}

    \begin{remark} \label{rmk:orientazione_opposta_1}
        Sia $(M, \Theta = \{ \Theta_x \}_{x \in M})$ una $m$-varietà orientata (con $\dim M > 1$ o $\partial M = \emptyset$).
        Allora si può definire l'\textbf{orientazione opposta} $-\Theta$:
        \[
            \boxed{-\Theta \defeq \{ -\Theta_x \}_{x \in M}.}
        \]
        In effetti $(M, -\Theta)$ è orientata: presa una parametrizzazione locale $g$ compatibile con
        $\Theta$, ristretta e traslata eventualmente a una palla di centro $0$, è sufficiente precomporla con una riflessione rispetto a un asse
        della palla per ottenere una parametrizzazione locale compatibile con $-\Theta$. \smallskip
    \end{remark}

    Questo ragionamento \underline{non} è attuabile sul bordo di una $1$-varietà:
    su $H^1$ una riflessione come quella sopracitata non è possibile. Questo ci suggerisce
    di modificare la definizione per il caso delle $1$-varietà bordate:

    \begin{definition}[$1$-varietà compatta orientata bordata]
        Sia $M$ una $1$-varietà connessa compatta con bordo (questo succede, per il Corollario \ref{cor:classificazione_dim_1}, se e solo se $M \cong [0, 1]$).
        Allora un'orientazione su $M$ è per definizione una famiglia
        $\Theta = \{\dif \varphi_t(\Theta_0)\}_{t \in [0, 1]}$ dove
        $\varphi : [0, 1] \to M$ è un diffeomorfismo. \smallskip

        Se $M$ è sconnessa, un'orientazione è un'orientazione su ciascuna componente connessa.
    \end{definition}

    \begin{remark} \label{rmk:orientazione_opposta_2}
        Se $M$ è una $1$-varietà connessa compatta con bordo, e $\Theta = \{\dif \varphi_t(\Theta_0)\}_{t \in [0, 1]}$ è
        una sua orientazione, allora
        \[
            \boxed{-\Theta \defeq \{-\dif \varphi_t(\Theta_0)\}_{t \in [0, 1]}}
        \]
        è una sua altra orientazione, indotta dalla precomposizione del diffeomorfismo $\varphi$ con
        una riflessione di $[0, 1]$ (e.g., $\psi(x) = 1-x$).
    \end{remark}

    \begin{proposition}
        Una varietà orientata e connessa, eventualmente con bordo, ammette esattamente due orientazioni.
    \end{proposition}

    \begin{proof}
        Poiché una varietà orientata ammette almeno due orientazioni (vd. Osservazione \ref{rmk:orientazione_opposta_1}
        e Osservazione \ref{rmk:orientazione_opposta_2}), per concludere la dimostrazione mostriamo che esistono al
        più due orientazioni. \smallskip

        Si fissi un'orientazione $\Theta$ per la varietà $M$. Sia $\Theta'$ un'altra orientazione.
        Dividiamo la dimostrazione in due casi:
        \begin{itemize}
            \item $\boxed{\dim M > 1 \text{ o } \partial M = \emptyset}$ Si definiscano i seguenti due insiemi:
                  \begin{equation*}
                      \begin{aligned}
                          A & \defeq \{ x \in M \mid \Theta_x = \Theta'_x \},  \\[1ex]
                          B & \defeq \{ x \in M \mid \Theta_x = -\Theta'_x \}.
                      \end{aligned}
                  \end{equation*}
                  Osserviamo che $A$ e $B$ sono disgiunti, e che la loro unione è la varietà $M$. Poiché
                  $M$ è connessa, mostrando che $A$ e $B$ sono aperti, necessariamente uno dei due deve essere
                  vuoto, da cui la tesi. \smallskip

                  Senza perdita di generalità, mostriamo solo che $A$ è aperto. Sia
                  $g : U \to g(U)$ una parametrizzazione locale di $x$ compatibile con $\Theta$ e $g(u) = x$, e sia $h : V \to h(V)$
                  una parametrizzazione compatibile con $\Theta'$ e $h(v) = x$. Assumiamo senza perdita di generalità che $g(U) = h(V)$.
                  \[\begin{tikzcd}
                          & M \\
                          U && V
                          \arrow["g", from=2-1, to=1-2]
                          \arrow["h"', from=2-3, to=1-2]
                          \arrow["{g\inv \circ h}", from=2-3, to=2-1]
                      \end{tikzcd}\]
                  Osserviamo che:
                  \[
                      \dif h_{v'} = \dif g_{g\inv(h(v'))} \circ \dif (g\inv \circ h)_{v'}.
                  \]
                  Poiché $g$ è compatibile con $\Theta$, si ha $\dif g(\Theta_0) = \Theta_x$, e così
                  per $h$ si ha $\dif h_x(\Theta_0) = \Theta'_x$. Quindi, per la precedente equazione:
                  \[
                      \dif g_u(\Theta_0) = \dif h_v(\Theta_0) \implies \dif (g\inv \circ h)_v (\Theta_0) = \Theta_0.
                  \]
                  Pertanto $\det(J(g\inv \circ h)_v) > 0$. Per continuità esiste allora un intorno $J$ di $v$ in cui:
                  \[
                      \det(J(g\inv \circ h)_{v'}) > 0, \quad \forall v' \in J.
                  \]
                  Questo si traduce nell'avere $\Theta_{h(v')} = \Theta'_{h(v')}$ su tutto $J$, e quindi
                  $h(J)$ è un aperto di $M$ contenente $x$ e contenuto in $A$; dunque $A$ è aperto.

            \item $\boxed{M \cong [0, 1]}$ Se $\varphi$ e $\psi$ sono due diffeomorfismi da $[0, 1]$ in $M$ che inducono
                  $\Theta$ e $\Theta'$, allora $\varphi\inv \circ \psi$ è un diffeomorfismo da $[0, 1]$ in sé. In quanto
                  tale, la sua derivata è ovunque non nulla, e il suo segno determina se $\Theta'$ è $\Theta$ o $-\Theta$.
        \end{itemize}
    \end{proof}

    \begin{definition}[Base positiva o negativa per $T_x M$]
        Sia $(M, \Theta)$ una varietà orientata. Allora una base per $T_x M$ si dice
        \textbf{positiva} se è della stessa orientazione di $\Theta_x$; altrimenti
        si dice \textbf{negativa}.
    \end{definition}

    \begin{definition}[Varietà di orientazione opposta]
        Data $(M, \Theta)$ una varietà orientata, indichiamo con $-M$ la varietà
        $(M, -\Theta)$, dove $-\Theta$ è l'unica altra orientazione possibile.
    \end{definition}

    \begin{remark}
        Una varietà è sempre ``localmente orientabile'': è sufficiente
        prendere l'orientazione indotta da un'unica parametrizzazione locale.
    \end{remark}

    \subsection{Orientabilità di \texorpdfstring{$m$}{m}-varietà immerse in \texorpdfstring{$\RR^m$}{ℝᵐ}}

    \begin{proposition}[$m$-varietà immerse in $\RR^m$ sono orientabili] \label{prop:orientazione_immersa_Rm}
        Sia $M$ una $m$-varietà immersa in $\RR^m$. Allora
        $M$ è orientabile secondo l'orientazione canonica di $\RR^m$.
    \end{proposition}

    \begin{proof}
        Sia $x \in M$. Osserviamo
        che deve valere necessariamente $T_x M = \RR^m$, dal momento che
        $T_x M$ è uno spazio vettoriale $m$-dimensionale immerso in $\RR^m$.
        Definiamo allora $\Theta_x \defeq \Theta_0$,
        dove $\Theta_0$ è l'orientazione canonica di $\RR^m$. \smallskip

        Se $g : U \subseteq H^m \to g(U)$ è una parametrizzazione locale di $x$, possiamo restringere e
        riparametrizzare $g$
        in modo tale che il suo dominio sia una semipalla di $\RR^m$. Se $\dif g_{-}$ preserva
        l'orientazione canonica, scegliamo $g$ come parametrizzazione locale compatibile
        per $\Theta \defeq \{\Theta_x\}_{x \in M}$. Altrimenti possiamo precomporre $g$
        con una riflessione rispetto all'asse della semipalla e ottenere una nuova parametrizzazione,
        stavolta compatibile. \smallskip

        Infatti, $g$ si estende a un diffeomorfismo tra aperti di $\RR^m$ e,
        per il Teorema della permanenza del segno, lo jacobiano deve essere localmente o positivo
        o negativo, e quindi $g$ preserva localmente l'orientazione.
        Dunque $M$ è orientabile secondo l'orientazione canonica di $\RR^m$.
    \end{proof}

    \subsection{Orientazione nel prodotto di due varietà orientate}

    \begin{definition}[Orientazione prodotto per varietà]
        Siano $(M, \Theta^M)$ e $(N, \Theta^N)$ due varietà orientate.
        Si definisce l'\textbf{orientazione prodotto} su $M \times N$
        come l'orientazione $\Theta^{M \times N}$ tale per cui:
        \[
            \boxed{\Theta_{(x, y)}^{M \times N} = \Theta_x^M \times \Theta_y^N, \quad \forall x \in M, y \in N}
        \]
        dove $\Theta_x^M \times \Theta_y^N$ è l'orientazione prodotto su
        $T_{(x, y)} M \times N \cong T_x M \times T_y N$ indotta da
        $\Theta_x^M$ e $\Theta_y^N$. \smallskip

        Il prodotto di parametrizzazioni locali compatibili induce parametrizzazioni
        locali compatibili con l'orientazione prodotto appena definita.
    \end{definition}

    \subsection{Semispazio interno o esterno}

    \begin{remark}[Il semispazio interno è ben definito]
        L'orientazione locale di una $m$-varietà $M$ con bordo determina sempre la scelta di uno dei semispazi di
        $T_x M \setminus T_x \partial M$ per ogni $x$ sul bordo $\partial M$. \smallskip

        Se infatti $g : U \to g(U)$ è una parametrizzazione di un punto $x \in \partial M$ con $g(u) = x$,
        il semispazio scelto è proprio:
        \[ \dif g_u (H^m \setminus \partial H^m). \]
        Questo semispazio non dipende dalla parametrizzazione scelta. Se infatti $h : V \to h(V)$ è
        un'altra parametrizzazione con $h(v) = x$, a meno di restringere le mappe possiamo considerare
        la funzione di transizione $h\inv \circ g$. Osserviamo che:
        \[
            \dif g_u = \dif h_v \circ \dif (h\inv \circ g)_u.
        \]
        Mostrando allora che $\dif (h\inv \circ g)_u(H^m \setminus \partial H^m) = H^m \setminus \partial H^m$,
        otteniamo la tesi. \smallskip

        Osserviamo che $h\inv \circ g$ è una funzione da un aperto intersecante il bordo di $H^m$ in
        un altro aperto dello stesso tipo. Pertanto $\dif (h\inv \circ g) : \RR^{m-1} \times \RR \to \RR^{m-1} \times \RR$
        deve mandare lo spazio tangente $T_u \partial U = \partial H^m$ in $T_v \partial V = \partial H^m$, dal momento che $h\inv \circ g$
        si restringe a una parametrizzazione di $\partial V$ a partire da $\partial U$. \smallskip

        Quindi $J(h\inv \circ g)_u$ deve essere della seguente forma:
        \[
            J(h\inv \circ g)_u = \begin{pmatrix}
                * & \vline & *       \\
                \hline
                0 & \vline & \lambda
            \end{pmatrix}.
        \]
        Se $h\inv \circ g = (\Psi_1, \Psi_2)$, con $\Psi_1$ funzione a valori in $\RR^{m-1}$ e
        $\Psi_2$ a valori in $\RR$, allora:
        \[
            \lambda = \pd{\Psi_2}{x_m}(u) = \lim_{\eps \to 0^+} \frac{\overbrace{\Psi_2(x + \eps x_m)}^{\mathclap{\geq \, 0}} - \overbrace{\Psi_2(x)}^{\mathclap{= \, 0}}}{\eps} \geq 0.
        \]
        Quindi $\dif (h\inv \circ g)_u(H^m) = H^m$, da cui segue poi facilmente la tesi.
    \end{remark}

    \begin{definition}
        Data una $m$-varietà bordata $M$ e un punto $x$ appartenente al bordo
        $\partial M$, si definisce il \textbf{semispazio interno} riferito a
        $x$ come:
        \[
            \boxed{\dif g_u(H^m \setminus \partial H^m),}
        \]
        dove $g$ è una parametrizzazione locale di $x$ con $g(u) = x$. I suoi
        vettori sono detti \textbf{interni}. \smallskip

        Si dice \textbf{semispazio esterno} il semispazio complementare a quello
        interno rispetto al taglio dell'iperpiano $T_x \partial M$ in
        $T_x M$, e i suoi vettori sono detti \textbf{esterni}.
    \end{definition}

    \begin{lemma} \label{lem:0_interno_1_esterno}
        Sia $M \cong [0, 1]$. Si fissi un diffeomorfismo
        $\varphi : [0, 1] \to M$. Allora
        $\dif \varphi_0(1)$ è un vettore interno,
        mentre $\dif \varphi_1(1)$ è esterno.
    \end{lemma}

    \begin{proof}
        Si fissi $\eps$ tale per cui $0 < \eps < 1$. Poniamo
        $I_0 = [0, \eps)$ e $I_1 = (\eps, 1]$.
        Osserviamo che $\restr{\varphi}{I_0}$ è una
        parametrizzazione locale di $\varphi(0)$, una
        volta identificato $I_1$ naturalmente con $H^1$.

        Allora $\dif \varphi_0(1)$ è interno dal momento
        che $\dif \varphi_0\inv(\dif \varphi_0(1)) = 1 > 0$. \smallskip

        Per quanto riguarda invece $\varphi(1)$, una parametrizzazione
        locale è data su $I_1$ con $H^1$
        secondo $\Psi(t) = \varphi(1-t)$. \smallskip

        Osserviamo che $\dif \Psi_0\inv (\dif \varphi_1(1)) = -1 < 0$, e dunque
        $\dif \varphi_1(1)$ è invece esterno.
    \end{proof}

    \subsection{Orientazione sul bordo della varietà}

    \begin{remark}[L'orientazione indotta sul bordo è ben definita -- \textit{esistenza}]
        Sia $x \in \partial M$ un punto della $m$-varietà orientata $(M, \Theta)$, con $m > 1$.
        Sia $g : U \subseteq H^m \to g(U)$ una parametrizzazione locale di $x$ con $g(u) = x$ che
        sia compatibile con l'orientazione $\Theta$. Allora:

        \begin{itemize}
            \item per definizione, $\dif g_u(-e_m)$ è
                  un vettore \textit{esterno} per $x$;
            \item $\{\dif g_u(e_i)\}_{i=1 -- m-1}$ è una base
                  di $T_x \partial M$, dacché $\restr{g}{\partial H^m}$
                  si identifica come una parametrizzazione locale di $x$
                  in $\partial M$;
            \item $\{\dif g_u(e_i)\}_{i=1 -- m}$ è una base positiva
                  di $T_x M$, dacché $g$ è compatibile e $\{e_i\}_i$ ha
                  l'orientazione canonica in $\RR^m$.
        \end{itemize}
    \end{remark}

    \begin{remark}[L'orientazione indotta sul bordo è ben definita -- \textit{unicità}]
        Sia $x \in \partial M$ un punto della $m$-varietà orientata $(M, \Theta)$, con $m > 1$.
        Sia $\{v_1, \ldots, v_n\}$ una base di $T_x M$ tale per cui
        $v_1$ un vettore \underline{esterno} per $x \in \partial M$ e
        $\{v_2, \ldots, v_n\}$ è base di $T_x \partial M$. \smallskip

        Sia $\{v_1', \ldots, v_n'\}$ un'altra tale base di $T_x M$.
        Sia $g : U \to g(U)$ una parametrizzazione locale di $x$ con
        $g(u) = x$. Allora $\dif g_u$ è un isomorfismo, e in quanto tale
        lascia invariate le relazioni di orientazioni delle basi di $T_u U = \RR^m$
        quando portate in $T_x M$. \smallskip

        Sia $w_i = \dif g_u\inv(v_i)$ e sia $w_i' = \dif g_u\inv(v_i')$. Dal momento
        che $v_1$ e $v_1'$ sono vettori esterni, si deve avere necessariamente
        $(w_1)_m$, $(w_1')_m < 0$. Dal momento che $\restr{g}{\partial H^m}$ si identifica
        naturalmente come una parametrizzazione locale di $x$ in $\partial M$, e che i $v_i$ e i $v_i'$ per $i > 1$ formano
        una base di $T_x \partial M$, si ha $(w_i)_m = (w_i')_m = 0$ per ogni $i > 1$. \smallskip

        Dunque la matrice di cambio di base da $\{v_1, \ldots, v_n\}$ a
        $\{v_1', \ldots, v_n'\}$ è della seguente forma:
        \[
            M = \begin{pmatrix}
                \lambda & \vline & 0 \\
                \hline
                *       & \vline & A
            \end{pmatrix},
        \]
        dove $\lambda > 0$ affinché $w_1$ e $w_1'$ abbiano ancora lo stesso segno
        sull'ultima coordinata. Dal momento che $\{v_1, \ldots, v_n\}$ e
        $\{v_1', \ldots, v_n'\}$ sono basi positive di $T_x M$, allora hanno
        stesso orientazione, e quindi $\det(M) > 0$. Ne segue che $\det(A) > 0$. \smallskip

        Osserviamo che $A$ è proprio la matrice di cambio di base da $\{v_2, \ldots, v_n\}$ a
        $\{v_2', \ldots, v_n'\}$. Dunque queste due basi hanno stessa orientazione.
    \end{remark}

    \begin{remark}[L'orientazione indotta sul bordo è ben definita -- \textit{è effettivamente un'orientazione}]
        Sia $x \in \partial M$ un punto della $m$-varietà orientata $(M, \Theta)$, con $m > 1$.
        Denotiamo con $\Theta_x^{\partial M}$ l'orientazione indotta da $\{v_2, \ldots, v_n\}$
        su $T_x \partial M$ da una base positiva $\{v_1, v_2, \ldots, v_n\}$ di $T_x M$
        con $v_1$ \underline{esterno} e $\{v_2, \ldots, v_n\}$ base di $T_x \partial M$. \smallskip

        Definiamo:
        \[ \Theta^{\partial M} \defeq \{ \Theta_x^{\partial M} \}_{x \in \partial M}. \]
        Data $g : U \subseteq H^m \to g(U) \subseteq M$ parametrizzazione locale
        compatibile di
        $x \in \partial M$ in $M$, $\restr{g}{\partial U}$, identificata come parametrizzazione da $\RR^{m-1}$,
        è compatibile rispetto a $\Theta^{\partial M}$, a meno di restringimento del dominio a una palla
        con conseguente riflessione rispetto a un asse. \smallskip

        Si può infatti estendere in tal caso la base canonica di $\RR^{m-1} \cong \partial H^m$ a una base di $\RR^m$ con l'aggiunta di
        un vettore la cui immagine tramite $\dif g_{-}$ risulta essere sempre esterna. Quindi
        $\Theta^{\partial M}$ è un'orientazione per $\partial M$.
    \end{remark}

    \begin{definition}[Orientazione indotta sul bordo]
        Sia $(M, \Theta)$ una $m$-varietà bordata e orientata con $m > 1$. Per $x \in \partial M$ denotiamo con
        $\Theta_x^{\partial M}$ l'orientazione indotta da $\{v_2, \ldots, v_n\}$ su $T_x \partial M$ da una base positiva $\{v_1, v_2, \ldots, v_n\}$ di $T_x M$
        con $v_1$ \underline{esterno} e $\{v_2, \ldots, v_n\}$ base di $T_x \partial M$. \smallskip

        Si definisce allora $\Theta^{\partial M} = \{ \Theta_x^{\partial M} \}_{x \in \partial M}$ come
        l'\textbf{orientazione indotta sul bordo} (o \textit{orientazione di bordo}) per $\partial M$.

        Per $M \cong [0, 1]$ orientata tramite un diffeomorfismo $\varphi : [0, 1] \to M$, si associa
        $-1$ a $\varphi(0)$ e $+1$ a $\varphi(1)$.
    \end{definition}

    \begin{corollary} \label{cor:bordo_orientabile_immersa_Rm}
        Sia $M$ una $m$-varietà orientabile immersa in $\RR^m$. Allora
        $\partial M$ è orientabile con l'orientazione indotta sul bordo dall'orientazione
        canonica di $\RR^m$.
    \end{corollary}

    \begin{proof}
        Segue immediatamente dalla Proposizione \ref{prop:orientazione_immersa_Rm}.
    \end{proof}

    \begin{corollary}
        $S^n$ è orientabile per ogni $n$.
    \end{corollary}

    \begin{proof}
        Segue immediatamente dal Corollario \ref{cor:bordo_orientabile_immersa_Rm},
        dacché $S^n = \partial D^{n+1}$ e $D^{n+1}$ è una $(n+1)$-varietà immersa
        in $\RR^{n+1}$.
    \end{proof}

    \section{Teoria del grado su \texorpdfstring{$\ZZ$}{ℤ}}

    \subsection{Grado intero rispetto a un valore regolare}

    \begin{definition}[Segno di un differenziale in un punto]
        Sia $M$ una varietà orientata con $\Theta^M$ e $N$ orientata
        con $\Theta^N$ e $\dim M = \dim N$. Se $f : M \to N$ è una mappa liscia, si definisce
        il \textbf{segno di $\dif f_x$} per un punto regolare $x \in M$ come:
        \[
            \boxed{\sgn(\dif f_x) = \begin{cases}
                    +1 & \text{se } \dif f_x(\Theta^M) = \Theta^N,  \\
                    -1 & \text{se } \dif f_x(\Theta^M) = -\Theta^N.
                \end{cases}}
        \]
    \end{definition}

    \begin{definition}[Grado intero di un valore regolare]
        Sia $M$ una varietà chiusa e sia $N$ una varietà connessa con $\partial N = \emptyset$
        e $\dim M = \dim N$. Sia $M$ orientata con $\Theta^M$ e sia $N$ orientata con $\Theta^N$. \smallskip

        Se $y \in N$ è regolare per $f : M \to N$ liscia, si definisce il \textbf{grado di $f$ rispetto a $y$} come:
        \[
            \boxed{\deg(f; y) \defeq \sum_{x \in f\inv(y)} \sgn(\dif f_x).}
        \]
    \end{definition}

    \begin{remark}[Il grado intero è localmente costante] \label{rmk:grado_loc_costante}
        Se $y \in N$ è regolare, possiamo scegliere per il Lemma \ref{lem:pila} un intorno $I_1$
        sul quale $\abs{f\inv(-)}$ è costante. \smallskip

        Inoltre, per il Teorema della permanenza del
        segno applicato sul determinante dello jacobiano di $h\inv \circ f \circ g$, dove
        $g : U \to g(U)$ parametrizza localmente $x \in f\inv(y)$ e $h : V \to h(V)$
        un intorno di $y$, esiste un intorno $I_2$ di $y$ sul quale le orientazioni sono
        preservate allo stesso modo in cui lo sono preservate dalle controimmagini di $y$. \smallskip

        Poiché i valori regolari sono aperti (vd. Lemma \ref{lem:valori_regolari_aperto}),
        possiamo allora prendere un intorno aperto di valori regolari
        in $I_1 \cap I_2$ entro cui $\deg(f; -)$ è costante; quindi il grado intero è \textit{localmente}
        costante.
    \end{remark}

    \subsection{Grado di una mappa estendibile dal bordo}

    \begin{lemma}[Il grado di una mappa estendibile dal bordo è nullo] \label{lem:grado_mappa_estendibile}
        Sia $X$ una varietà compatta e orientata con bordo non nullo. Sia $N$
        una varietà connessa, orientata, senza bordo e con $\dim X = \dim N + 1$.
        Sia $F : X \to N$ una mappa liscia. \smallskip

        Se $y \in N$ è un valore regolare per $f \defeq \restr{F}{\partial X}$, allora
        $\deg(f; y) = 0$.
    \end{lemma}

    \begin{proof}
        Denotiamo con $n$ la dimensione di $X$. \smallskip

        Grazie all'Osservazione \ref{rmk:grado_loc_costante}, sappiamo che esiste un intorno
        di $y$ fatto di valori regolari di $f$ su cui $\deg(f; -)$ è costante. Possiamo prendere allora in questo intorno,
        per il Teorema di Brown (per varietà, Corollario \ref{cor:brown}), un valore regolare
        $z$ di $F$. Per definizione, $z$ è valore regolare sia di $F$ che di $f$; mostrando che
        $\deg(f; z) = 0$, mostriamo dunque anche che $\deg(f; y) = 0$, per costruzione. \smallskip

        Per il Teorema \ref{thm:varietà_bordata_da_valore_regolare}, allora $F\inv(z)$ è una $1$-varietà
        compatta e chiusa topologicamente con bordo $F\inv(z) \cap \partial M = f\inv(z)$. Per il
        Corollario \ref{cor:classificazione_dim_1}, $F\inv(z)$ è unione di archi (con $2$ punti sul bordo ciascuno)
        e cerchi (che non hanno bordo). Per mostrare la tesi è sufficiente mostrare che per ogni arco
        i punti di bordo abbiano segno opposto sul differenziale, in modo tale che la somma complessiva dei
        segni sia ancora nulla. \smallskip

        Sia $A \subseteq F\inv(z) \subseteq X$ un tale arco. Fissiamo un diffeomorfismo $\varphi : [0, 1] \to A$.
        Poniamo $v_1(t) \defeq \dif \varphi_t (1)$. Allora vale:
        \[
            T_{\varphi(t)} X = \underbrace{T_{\varphi(t)} A}_{\mathclap{=\, \Span(v_1(t))}} \oplus (T_{\varphi(t)} A)^\perp,
        \]
        dove ricordiamo che, per una generalizzazione della Proposizione \ref{prop:tangente_valore_regolare},
        $\restr{\dif F_{\varphi(t)}}{(T_{\varphi(t)} A)^\perp} : (T_{\varphi(t)} A)^\perp \to T_z N$ è un isomorfismo. \smallskip

        Fissiamo una base positiva $\{w_2, \ldots, w_n\}$ di $T_y N$. Possiamo allora porre:
        \[
            v_i(t) = \restr{\dif F_{\varphi(t)}}{(T_{\varphi(t)} A)^\perp}\inv(w_i), \quad i = 2, \ldots, n.
        \]
        Poiché $v_1(t)$ è ortogonale agli altri $v_i(t)$, $\{v_i(t)\}_i$ è una base
        di $T_{\varphi(t)} X$. \smallskip

        Se poniamo $A \defeq \{t \mid \{v_i(t)\}_i \text{ positiva}\} \subseteq [0, 1]$,
        allora $A$ è un aperto di $[0, 1]$. Infatti, per il Teorema della permanenza del segno applicato
        allo jacobiano di una parametrizzazione locale, in un
        intorno di $t_0$ si mantiene la stessa orientazione.
        Analogamente anche il complementare
        di $A$ è aperto. Dacché $[0, 1]$ è connesso, allora uno tra $A$ e il suo complementare è vuoto: in altre
        parole $\{v_i(t)\}_i$ è sempre positiva per ogni $t \in [0, 1]$, o altrimenti è sempre negativa. \smallskip

        A meno di invertire l'orientazione di $\varphi$ usando al suo posto $\varphi(1 - t)$, possiamo
        assumere che $\basis(t) \defeq \{v_i(t)\}_i$ sia sempre positiva. \smallskip

        Mostriamo che $\sgn(\dif f_{\varphi(0)}) = -1$. Sia $\basis' \defeq \{v_1(0), v_2'(0), \ldots, v_n'(0)\}$
        una base di $T_{\varphi(0)} \partial X$ tale per cui $\{ \dif F_{\varphi(0)} (v_i'(0)) \}_{i \geq 2} $ è
        una base positiva di $T_z N$. Se mostriamo che $\basis$ e $\basis'$ hanno la stessa orientazione,
        dal momento che $\basis$ è positiva e $v_1(0)$ è \underline{interno} per il Lemma \ref{lem:0_interno_1_esterno},
        si deduce dalla definizione di orientazione di bordo che
        $\{ v_2'(0), \ldots, v_n'(0)\}$ è negativa; in particolare $\dif f_{\varphi(0)}$ invertirebbe l'orientazione,
        e quindi avrebbe segno $-1$. \smallskip

        Consideriamo la matrice di cambio di base da $\basis$ a $\basis'$, della seguente forma:
        \[
            M \defeq \begin{pmatrix}
                1 & \vline & * \\
                \hline
                0 & \vline & B
            \end{pmatrix}.
        \]
        Osserviamo che $B$ è la matrice di cambio di base da $\{ \dif F_{\varphi(0)} (v_i(0)) \}_{i \geq 2}$ a $\{ \dif F_{\varphi(0)} (v_i'(0)) \}_{i \geq 2}$.
        Per costruzione, queste tali basi hanno la stessa orientazione; dunque
        $\det(B) > 0$, da cui si deduce che $\det(M) > 0$, e quindi che $\basis$ e
        $\basis'$ hanno effettivamente la stessa orientazione. \smallskip

        Per $\varphi(1)$ il ragionamento è del tutto analogo, ma essendo $v_1(1)$ \underline{esterno} per
        il Lemma \ref{lem:0_interno_1_esterno}, il segno sarà invece $+1$. Questo conclude la dimostrazione
        per le osservazioni fatte in precedenza.
    \end{proof}

    \subsection{Passaggio per omotopia e buona definizione del grado intero di una mappa}

    \begin{lemma}[di omotopia, per il grado intero] \label{lem:omotopia_intero}
        Siano $M$ e $N$ due varietà orientate con $M$ chiusa e $\dim M = \dim N$.
        Sia $F : M \times [0, 1] \to N$ un'omotopia $C^\infty$ con $f \defeq F_0$
        e $g \defeq F_1$. Se $y \in N$ è un valore regolare comune a $f$ e $g$,
        allora:
        \[ \boxed{\deg(f; y) = \deg(g; y).} \]
    \end{lemma}

    \begin{proof}
        Dal momento che $M \times [0, 1]$ è un prodotto, allora acquisisce l'orientazione
        prodotto di quella di $M$ con quella canonica di $[0, 1]$. Grazie
        all'Osservazione \ref{rmk:grado_loc_costante} e al Teorema di Brown (per varietà, Corollario \ref{cor:brown}),
        possiamo assumere senza perdita di generalità che $y$ sia un valore regolare anche di $F$. \smallskip

        Sia $\basis_x$ una base positiva di $T_x M$ per $x \in M$. Allora,
        per l'Osservazione \ref{rmk:regola_segni} la base $\basis_x \cup \{1\}$
        è positiva per $T_{(x, t)} M \times [0, 1] \cong T_x M \times T_t [0, 1]$. \smallskip

        Ricordiamo che il bordo di $M \times [0, 1]$ è $M \times \{0\} \sqcup M \times \{1\}$
        per \ref{rmk:m_01_varietà}. Studiamo l'orientazione indotta su tale bordo.
        Poiché $(0, 1) \in T_{(x, 0)} M \times [0, 1]$ è interno per il Lemma \ref{lem:0_interno_1_esterno},
        $\{-\basis_x\} \cup \{-1\}$ è positiva per $(x, 0)$, e quindi $\basis_x$ è una base negativa di $(x, 0)$ sul
        bordo. Dunque $M \times \{0\} \cong M$ è orientata come $-M$. Analogamente, $(0, 1) \in T_{(x, 1)} M \times [0, 1]$ è
        esterno, quindi $M \times \{1\} \cong M$ è orientata come $M$. \smallskip

        Allora, per il Lemma \ref{lem:grado_mappa_estendibile}, si ha:
        \[
            \deg(g; y) - \deg(f; y) = \deg(\restr{F}{\partial(M \times [0, 1])}; y) = 0,
        \]
        da cui la tesi.
    \end{proof}

    \begin{theorem}[Il grado intero è ben definito] \label{thm:grado_intero_ben_definito}
        Siano $M$ e $N$ varietà con $M$ chiusa, $N$ connessa e $\dim M = \dim N$. Se
        $f : M \to N$ è liscia e $y$, $z \in N$ sono suoi valori regolari, allora:
        \[
            \boxed{\deg(f; y) = \deg(f; z).}
        \]
    \end{theorem}

    \begin{proof}
        Per il Lemma \ref{lem:omogeneità}, esiste un diffeomorfismo $h : N \to N$
        con $h(y) = z$ isotopo all'identità $\id_N$. Sia $H : N \times [0, 1] \to N$
        una tale isotopia. \smallskip

        Consideriamo la mappa $f(t) = \sgn(\dif (H_t)_y (\Theta^N_y))$. Poiché
        $H_1 = \id_N$, si ha chiaramente $f(1) = 1$. Inoltre $f$ è continua e localmente
        costante, utilizzando l'usuale argomento sulla permanenza del segno
        sul determinante dello jacobiano, una volta scelta una parametrizzazione locale.
        Dunque, poiché $[0, 1]$ è connesso, $f$ deve essere costantemente uguale a $1$.
        Quindi $\dif (H_t)_y (\Theta^N_y) = \Theta^N_{H_t(y)}$ per ogni $t \in [0, 1]$. \smallskip

        Segue quindi che $\dif h_y (\Theta^N_y) = \Theta^N_z$, da cui per la regola della catena:
        \[
            \deg(h \circ f; z) = \deg(f; y).
        \]
        D'altra parte $h \circ f$ e $f$ sono $C^\infty$-omotope tramite $H \circ (f \times \id_{0, 1})$,
        e quindi, per il Lemma \ref{lem:omotopia_intero}:
        \[
            \deg(f; y) = \deg(h \circ f; z) = \deg(f; z).
        \]
    \end{proof}

    \begin{definition}[Grado intero di una mappa liscia]
        Sia $M$ una varietà chiusa e sia $N$ una varietà connessa.
        Siano $M$ e $N$ della stessa dimensione e orientate. Se
        $f : M \to N$ è una mappa liscia, si definisce il suo
        \textbf{grado intero} come:
        \[
            \boxed{\deg(f) \defeq \deg(f; y),}
        \]
        dove $y$ è un valore regolare qualsiasi di $f$.
    \end{definition}

    \begin{theorem} \label{thm:fondamentale_grado_intero}
        Sia $M$ una varietà chiusa e sia $N$ una varietà connessa.
        Siano $M$ e $N$ della stessa dimensione e orientate.
        Se $f$ e $g$ sono due mappe lisce $C^\infty$-omotope da $M$ in $N$,
        allora:
        \[
            \boxed{\deg(f) = \deg(g).}
        \]
    \end{theorem}

    \begin{proof}
        Per il Lemma \ref{lem:valori_regolari_aperto}, i valori regolari di $f$ formano un aperto in $N$. Allora, per il Teorema
        di Brown (per varietà, Corollario \ref{cor:brown}), esiste in questo
        aperto anche un valore regolare di $g$. Per il Lemma \ref{lem:omotopia_intero},
        dunque $f$ e $g$ condividono lo stesso grado intero.
    \end{proof}

    \begin{corollary}
        Sia $M$ una varietà chiusa e connessa. Se $f : M \to M$
        è un diffeomorfismo di grado $\deg(f) = -1$, allora
        $f$ non è omotopa all'identità $\id_M$, né a una mappa costante $c_x$ per
        $x \in M$.
    \end{corollary}

    \begin{proof}
        La tesi è un'immediata conseguenza del Teorema \ref{thm:fondamentale_grado_intero},
        dal momento che $\deg(\id_M) = 1$ e $\deg(c_x) = 0$ (infatti $c_x$ non è surgettiva).
    \end{proof}

    \begin{proposition}[Il grado è moltiplicativo] \label{prop:grado_moltiplicativo}
        Sia $M$ una varietà orientata, chiusa e connessa. Se $f$ e $g$ sono due
        mappe lisce da $M$ in sé stessa, allora:
        \[
            \boxed{\deg(f \circ g) = \deg(f) \deg(g).}
        \]
    \end{proposition}

    \begin{proof}
        Sia $z$ un valore regolare di $f \circ g$. Sia $x \in (f \circ g)\inv(z)$.
        Allora, per la regola della catena:
        \[
            \dif (f \circ g)_x = \dif f_{g(x)} \circ \dif g_x.
        \]
        Dal momento che $z$ è regolare, $x$ è un punto regolare e dunque
        $\dif (f \circ g)_x$ è un isomorfismo. Dunque necessariamente
        $\dif f_{g(x)}$ e $\dif g_x$ sono isomorfismi, ovverosia
        $x$ è regolare anche per $g$, mentre $g(x)$ lo è per $f$. \smallskip

        Sia $f\inv(z) = \{y_1, \ldots, y_n\}$.
        Da quanto appena visto, ogni $y_i$ è valore regolare.
        Allora:
        \[
            g\inv(f\inv(z)) = \bigsqcup_{i = 1}^n g\inv(y_i).
        \]

        Allora, sfruttando anche l'Osservazione \ref{rmk:regola_segni}:
        \[
            \deg(f \circ g) = \sum_{x \in g\inv(f\inv(z))} \sgn(\dif f_{g(x)}) \sgn(\dif g_x).
        \]
        Raccogliendo i termini utilizzando $f\inv(z)$ si ottiene dunque, anche usando il Teorema \ref{thm:grado_intero_ben_definito}:
        \[
            \deg(f \circ g) = \sum_{y_i \in f\inv(y_i)} \sgn(\dif f_{y_i}) \deg(g) = \deg(f) \deg(g).
        \]
    \end{proof}

    \subsection{Grado di \texorpdfstring{$z^k$}{zᵏ}, delle riflessioni e della mappa antipodale su \texorpdfstring{$S^1$}{S¹}}

    \begin{lemma} \label{lem:grado_zk}
        Sia $f_k : S^1 \to S^1$ tale per cui:
        \[
            f_k(z) = z^k \in \CC,
        \]
        dove si è identificato $S^1 \subseteq \RR^2$ in $\CC$.
        Allora $1$ è un valore regolare di $f_k$ e $\deg(f_k) = k$. \smallskip

        Quindi, per $k \neq 0$, $f_k$ \underline{non} può estendersi a una mappa
        liscia da $D^2$ a $S^1$.
    \end{lemma}

    \begin{proof}
        Consideriamo l'elemento $1 \in S^1$. Osserviamo che
        $f_k\inv(1)$ è l'insieme delle radici $k$-esime dell'unità,
        e quindi contiene esattamente $k$ elementi (a meno del segno). \smallskip

        La funzione $f_k$ si estende a $F_k(z) = z^k$ su tutto $\CC$. Osserviamo
        che $T_1 S^1 = \Span(i)$. Per determinare allora il segno di $\dif (f_k)_1$ è
        sufficiente considerare la seguente derivata:
        \[
            \dif (f_k)_1 (i) = \dertime{e^{i t k}}{t=0} = k i.
        \]
        Dunque $\dif (f_k)_1$ è un isomorfismo per $k \neq 0$, e preserva l'orientazione
        se $k > 0$, mentre non la preserva se $k < 0$. \smallskip

        Sia ora $\xi = e^{i \theta_0} \in f_k\inv(1)$. Consideriamo il diffeomorfismo $h_\xi : S^1 \to S^1$ tale
        per cui:
        \[
            h_\xi(z) = \xi \cdot z,
        \]
        ovverosia la rotazione indotta da $\xi$. Si verifica facilmente che:
        \[
            H : S^1 \times [0, 1] \to S^1, \quad H(z, t) = e^{i \theta_0 t} z
        \]
        è un'omotopia liscia da $\id_{S^1}$ a $h_\xi$. Dunque $\deg(h_\xi; 1) = \deg(\id_{S^1}; 1) = 1$. \smallskip

        Osserviamo che $f_k = f_k \circ h_\xi$. Quindi:
        \[
            \dif (f_k)_1 = \dif (f_k \circ h_\xi)_1 = \dif (f_k)_\xi \circ \dif (h_\xi)_1 = \dif (f_k)_\xi \circ h_\xi.
        \]
        Dal momento che $h_\xi$ è invertibile, si deduce che $\dif (f_k)_\xi$ è sempre un isomorfismo, e
        dunque che $1$ è un valore regolare. Inoltre, dalla stessa uguaglianza si deduce per la
        moltiplicatività del segno dei differenziali che $\sgn(\dif (f_k)_1) = \sgn(\dif (f_k)_\xi)$. \smallskip

        Si conclude facilmente allora che $\deg(f) = \deg(f; 1) = k$. L'ultima affermazione è conseguenza
        del Lemma \ref{lem:grado_mappa_estendibile}.
    \end{proof}

    \begin{remark}[Un diffeomorfismo ha grado $1$ o $-1$] \label{rmk:grado_diffeomorfismo}
        Per un diffeomorfismo $f : M \to M$ su $M$ chiusa e connessa,
        possono esistere solo due gradi, $+1$ o $-1$, dacché l'insieme
        controimmagine di un valore regolare contiene un singolo elemento. \smallskip

        In particolare, $\deg(f) = 1$ se e solo se per un elemento $x \in M$, $\dif f_x$
        preserva l'orientazione.
    \end{remark}

    \begin{lemma} \label{lem:grado_riflessione}
        Sia $r_i : S^n \to S_n$ la riflessione tale per cui:
        \[
            r_i(x_1, \ldots, x_i, \ldots, x_{n+1}) \defeq (\underbrace{x_1, \ldots}_{\textnormal{invariati}}, -x_i, \underbrace{\ldots, x_{n+1}}_{\textnormal{invariati}}).
        \]
        Allora $\deg(r_i) = -1$.
    \end{lemma}

    \begin{proof}
        Per l'Osservazione \ref{rmk:grado_diffeomorfismo} è sufficiente studiare $\dif (r_i)_{e_i}$.
        Per il Corollario \ref{cor:sn_è_varietà},
        si ha:
        \[
            T_{e_i} S^n = e_i^\perp = T_{-e_i} S^n.
        \]
        $r_i$ si estende con la stessa formula a un diffeomorfismo $\tilde{r_i}$ su $D^{n+1}$. Poiché
        $\tilde{r_i}$ è lineare e $T_{e_i} D^{n+1} = \RR^{n+1}$, si ha
        $\dif \tilde{r_i}_{e_i} = \tilde{r_i}$. Tale differenziale manda la base $\{e_i, e_1, \ldots, e_{n+1}\}$
        in $\{-e_i, e_1, \ldots, e_{n+1}\}$. Queste due basi hanno orientazione diversa, e quindi solo
        una di queste induce l'orientazione canonica su $\RR^{n+1}$. \smallskip

        Osserviamo che $e_i$ è esterno per sé stesso; allo stesso modo $-e_i$ è esterno per sé stesso.
        Segue allora facilmente che $\{e_1, \ldots, e_{n+1}\}$ cambia orientazione tramite $r_i$,
        e quindi $\deg(r_i) = -1$.
    \end{proof}

    \begin{lemma} \label{lem:grado_antipodale}
        Sia $A : S^n \to S^n$ la mappa antipodale, ossia tale per cui $A(x) = -x$.
        Allora $\deg(A) = (-1)^{n+1}$.
    \end{lemma}

    \begin{proof}
        Segue immediatamente dalla Proposizione \ref{prop:grado_moltiplicativo} e il Lemma \ref{lem:grado_riflessione}.
    \end{proof}

    \subsection{Campi vettoriali tangenti su \texorpdfstring{$S^n$}{Sⁿ} e pettinabilità}

    \begin{definition}[Campo vettoriale tangente]
        Sia $M \subseteq \RR^k$ una varietà liscia, con o senza bordo. Un
        \textbf{campo vettoriale tangente} è una mappa liscia
        $v : M \to \RR^k$ tale per cui:
        \[
            \boxed{v(x) \in T_x M, \quad \forall x \in M.}
        \]
    \end{definition}

    \begin{definition}[Pettinabilità]
        Una varietà liscia, con o senza bordo, si dice \textbf{pettinabile}
        se ammette un campo vettoriale tangente mai nullo.
    \end{definition}

    \begin{theorem}[di pettinabilità della sfera] \label{thm:pettinabilità_sfera}
        $S^n$ è pettinabile se e solo se $n$ è dispari.
    \end{theorem}

    \begin{proof}
        Se $n$ è dispari, un campo vettoriale tangente mai nullo è il seguente:
        \[
            f(x_1, x_2, \ldots, x_n, x_{n+1}) = (-x_2, x_1, \ldots, -x_{n+1}, x_n) \in x^\perp.
        \]
        Sia ora $n$ pari.
        Supponiamo $v$ sia un campo vettoriale tangente mai nullo. Senza
        perdita di generalità, possiamo considerarlo unitario. Allora
        possiamo costruire un omotopia liscia dalla mappa
        antipodale $A$ a $\id_{S^n}$ nel seguente modo:
        \[
            H : S^n \times [0, 1] \to S^n, \quad H(x, t) = \cos(\pi t) x + \sin(\pi t) v(x).
        \]
        Tuttavia una tale omotopia \underline{non} può esistere per
        il Teorema \ref{thm:fondamentale_grado_intero}: l'identità
        ha grado $1$, mentre la mappa antipodale ha grado $(-1)^{n+1} = -1$ per
        il Lemma \ref{lem:grado_antipodale}. Quindi $v$ non può esistere per $n$ pari.
    \end{proof}

    \section{Indici di campi vettoriali su aperti di \texorpdfstring{$\RR^m$}{ℝᵐ}}

    \subsection{Zero isolato e indice di un campo in uno zero}

    \begin{definition}[Zero isolato]
        Sia $f : U \subseteq \RR^m \to \RR^m$ liscia con $U$ aperto. Allora
        $z$ si dice \textbf{zero isolato} di $f$ se esiste un raggio $\eps > 0$
        tale per cui $f$ in $B_\eps(z)$ ammette come unico zero $z$.
    \end{definition}

    \begin{remark}[L'indice è ben definito]
        Sia $\eps > 0$ tale per cui $z$ è unico zero per $f$ in $B_\eps(z)$.
        Sia $v_\eps : S^m \to \partial B_\eps(z)$ tale per cui:
        \[
            v_\eps(x) = z + \eps x.
        \]
        Osserviamo che $v_\eps$ preserva l'orientazione. Consideriamo
        \[
            \overline{f_\eps} \defeq \bigrestr{\frac{f}{\norm{f}}}{\partial B_\eps(z)}.
        \]
        Poiché $v_\eps$ preserva l'orientazione, $\deg(\overline{f_\eps}) = \deg(\overline{f_\eps} \circ v_\eps)$.
        Scelto un altro $\eps'$, possiamo definire un'omotopia $H$ nel seguente modo:
        \[
            H_t = \overline{f_{(1-t)\eps + \eps'}} \circ v_{(1-t)\eps + \eps'}.
        \]
        Allora, per il Teorema \ref{thm:fondamentale_grado_intero}:
        \[
            \deg(\overline{f_\eps}) = \deg(\overline{f_\eps} \circ v_\eps) = \deg(\overline{f_{\eps'}} \circ v_{\eps'}) = \deg(\overline{f_{\eps'}}).
        \]
    \end{remark}

    \begin{definition}[Indice di $f$ in $z$]
        Sia $f : U \subseteq \RR^m \to \RR^m$ liscia con $U$ aperto. Sia
        $z$ uno zero isolato di $f$. Si definisce allora l'\textbf{indice di $f$ in $z$}
        come:
        \[
            \boxed{\ind(f, z) \defeq \deg(\overline{f_\eps}), \quad \overline{f_\eps} \defeq \bigrestr{\frac{f}{\norm{f}}}{\partial B_\eps(z)},}
        \]
        dove $\eps$ è un raggio tale per cui $z$ è unico zero in $B_\eps(z)$.
    \end{definition}

    \begin{corollary}[Indice di $z^k$ in $0$]
        Sia $v_k : \CC \cong \RR^2 \to \CC$ tale per cui
        $v_k(z) = z^k$. Allora $\ind(v_k, 0) = k$.
    \end{corollary}

    \begin{proof}
        Segue immediatamente dal Lemma \ref{lem:grado_zk}.
    \end{proof}

    \subsection{Lemma di Hopf e teorema fondamentale dell'algebra}

    \begin{lemma}[di Hopf] \label{lem:hopf}
        Sia $X \subseteq \RR^m$ una $m$-varietà compatta con bordo. Sia
        $v : X \to \RR^m$ un campo vettoriale con zeri isolati e $\restr{v}{\partial X}$
        mai nullo. \smallskip

        Allora:
        \[
            \boxed{\sum_{z \in v\inv(0)} \ind(v, z) = \deg\left(\bigrestr{\frac{v}{\norm{v}}}{\partial X} : \partial X \to S^{m-1} \right).}
        \]
    \end{lemma}

    \begin{proof}
        Poiché gli zeri di $v$ sono isolati, possiamo togliere a $X$ una famiglia di dischi disgiunti $\{B_{\eps_i}(z_i)\}$ contenenti tali zeri.
        Allora, poiché tali dischi sono interni, l'orientazione indotta su $\partial W$ sarà:
        \[
            \partial W = \partial X \cup \bigsqcup_i -\partial B_{\eps_i}(z_i).
        \]
        Quindi, rispettando le orientazioni si ottiene:
        \begin{multline*}
            \deg\left(\bigrestr{\frac{v}{\norm{v}}}{\partial W} : \partial W \to S^{m-1} \right) = \\
            \deg\left(\bigrestr{\frac{v}{\norm{v}}}{\partial X} : \partial X \to S^{m-1} \right) - \sum_{z \in v\inv(0)} \ind(v, z).
        \end{multline*}
        Tuttavia, poiché $\bigrestr{\frac{v}{\norm{v}}}{\partial W}$ è la restrizione di $\frac{v}{\norm{v}} : W \to S^{m-1}$ sul bordo, per il Lemma
        \ref{lem:grado_mappa_estendibile}, si ha anche:
        \[
            \deg\left(\bigrestr{\frac{v}{\norm{v}}}{\partial W} : \partial W \to S^{m-1} \right) = 0,
        \]
        da cui segue immediatamente la tesi.
    \end{proof}

    \begin{theorem}[fondamentale dell'algebra]
        Sia $p(z) \in \CC[z]$ con $\deg(p) = n$. Allora:
        \[
            \boxed{\deg(p) = \sum_{z_0 \in p\inv(0)} \mult(p, z_0),}
        \]
        dove $\mult$ indica la molteplicità algebrica di uno zero in un polinomio.
    \end{theorem}

    \begin{proof}
        Dal momento che $p(x)$ non può avere più di $n$ zeri, questi sono sicuramente
        isolati e possiamo prendere inoltre una palla $B_r(0) \subseteq \CC$ con
        $p\inv(0) \subseteq B_r(0)$. Possiamo allora applicare il Lemma \ref{lem:hopf}
        su $p : \overline{B_r(0)} \to \CC$ e ottenere:
        \[
            \deg\left(\frac{p}{\norm{p}} : \partial B_r(0) \to S^1 \right) = \sum_{z_0 \in p\inv(0)} \ind(p, z_0).
        \]
        Mostriamo che il termine a sinistra coincide con $\deg(p)$ (1), e che $\ind(p, z_0)$ coincide
        con $\mult(p, z_0)$ (2), ottenendo infine la tesi.

        \begin{enumerate}[(1)]
            \item Supponiamo che $\deg(p)$ sia $n$ e che $p(z)$ sia dunque della seguente forma:
                  \[
                      p(z) = a_n z^n + \underbrace{a_{n-1} z^{n-1} + \ldots + a_0}_{\mathclap{g(z)}},
                  \]
                  dove $g(z) \defeq p(z) - a_n z^n$. \smallskip

                  Osserviamo che:
                  \begin{equation} \tag{*}
                      \lim_{\abs{z} \to \infty} \abs{\frac{g(z)}{z^n}} = 0.
                  \end{equation}
                  Una volta posto $p_t(z) = a_n z^n + t g(z)$, si ottiene:
                  \[
                      \abs{\frac{p_t(z)}{z^n}} \geq \abs{a_n} - t \abs{\frac{g(z)}{z^n}}.
                  \]
                  Allora, per (*), possiamo scegliere $r$ sufficientemente grande in modo tale che si verifichi sempre:
                  \[
                      \abs{\frac{p_t(z)}{z^n}} > 0, \quad z \in \partial B_r(0).
                  \]
                  In particolare, $p_t(z)$ non si annulla su $\partial B_r(0)$. Possiamo
                  allora considerare l'omotopia indotta da $\frac{p_t(z)}{\abs{p_t(z)}}$. Osserviamo che
                  $p_0(z) = a_n z^n$ e che $p_1(z) = p(z)$. Per il Teorema
                  \ref{thm:fondamentale_grado_intero}, si ha allora:
                  \[
                      \begin{split}
                          \deg\left( \frac{p(z)}{\abs{p(z)}} : \partial B_r(0) \to S^1 \right) = \hspace{2cm} \\
                          \hspace{2cm} \deg\left( \frac{a_n}{\abs{a_n}} \frac{z^n}{\abs{z^n}} : \partial B_r(0) \to S^1 \right),
                      \end{split}
                  \]
                  a cui, applicando il Lemma \ref{lem:grado_zk} e il fatto secondo cui la moltiplicazione per una costante di
                  fase è isotopa all'identità (vd. dimostrazione del Lemma \ref{lem:grado_zk}), si ottiene facilmente che:
                  \[
                      \deg\left( \frac{p(z)}{\abs{p(z)}} : \partial B_r(0) \to S^1 \right) = n = \deg(p).
                  \]

            \item Sia $z_0$ uno zero di $p(z)$. Allora $p(z)$ si scrive come:
                  \[
                      p(z) = (z - z_0)^\ell q(z),
                  \]
                  per un qualche polinomio $q(z) \in \CC[z]$, dove $\ell = \mult(p, z_0)$. Entro
                  una certa palla di raggio $\eps$ centrata in $z_0$, $q(z)$ non ha alcuno zero. Se
                  consideriamo la mappa $f : S^1 \to \partial B_{z_0}(\eps)$ tale per cui
                  $f(z) = z_0 + \eps z$, che preserva l'orientazione, allora si ha:
                  \[
                      \ind(p, z_0) = \deg\left(\frac{p \circ f}{\norm{p \circ f}} : \partial S^1 \to S^1 \right).
                  \]
                  Osserviamo che:
                  \[
                      \frac{p \circ f}{\norm{p \circ f}}\big(z\big) = \frac{z^\ell q(z_0 + \eps z)}{\abs{q(z_0 + \eps z)}}.
                  \]
                  Possiamo definire un'omotopia $H : S^1 \times [0, 1] \to S^1$ tale per cui:
                  \[
                      H_t(z) = \frac{z^\ell q(z_0 + t \eps z)}{\abs{q(z_0 + t \eps z)}},
                  \]
                  che porta $z^\ell \frac{q(z_0)}{\abs{q(z_0)}}$ in $\frac{p \circ f}{\norm{p \circ f}}$. Quindi,
                  per il Teorema \ref{thm:fondamentale_grado_intero}, si ha:
                  \[
                      \begin{split}
                          \deg\left(z^\ell \frac{q(z_0)}{\abs{q(z_0)}} : S^1 \to S^1\right) = \hspace{2cm} \\
                          \hspace{2cm} \deg\left(\frac{p \circ f}{\norm{p \circ f}} : \partial S^1 \to S^1 \right).
                      \end{split}
                  \]
                  Come visto nella dimostrazione del Lemma \ref{lem:grado_zk}, la moltiplicazione per elemento di $S^1$ è isotopa all'identità, e dunque:
                  \[
                      \deg\left(z^\ell \frac{q(z_0)}{\abs{q(z_0)}} : S^1 \to S^1\right) = \deg(z^\ell : S^1 \to S^1) = \ell,
                  \]
                  da cui segue, combinando i pezzi, che:
                  \[
                      \ind(p, z_0) = \ell.
                  \]
        \end{enumerate}
    \end{proof}

    \section{Campi vettoriali su varietà}

    \subsection{Indice di un campo vettoriale tangente su una varietà}

    \begin{fact}
        Sia $v : M \to \RR^k$ un campo vettoriale tangente
        della varietà $M \subseteq \RR^m$, con $z \in M$ zero isolato di $v$. \smallskip

        Se $g : U \subseteq \RR^m \to g(U)$ è una parametrizzazione di $M$ in $z$
        con $g(u) = z$,
        allora possiamo considerare il campo vettoriale
        $\xi : U \to \RR^m$ tale per cui:
        \[
            \xi(u) = (\dif g_u)\inv (v(g(u))), \quad \forall u \in U.
        \]
        Allora $\ind(\xi, g\inv(z))$ \underline{non} dipende dalla scelta
        della parametrizzazione scelta $g$.
    \end{fact}

    \begin{definition}[Indice di un campo vettoriale tangente su varietà rispetto a un punto]
        Sia $v : M \to \RR^k$ un campo vettoriale tangente
        della varietà $M \subseteq \RR^m$, con $z \in M$ zero isolato di $v$. \smallskip

        Si definisce allora l'\textbf{indice di $v$ in $z$} come:
        \[
            \boxed{\ind(v, z) \defeq \ind(\xi, g\inv(z))},
        \]
        dove $g : U \subseteq \RR^m \to g(U)$ è una parametrizzazione locale di $M$ in $z$
        con $g(u) = z$ e $\xi : U \to \RR^m$ è tale per cui:
        \[
            \boxed{\xi(u) = (\dif g_u)\inv (v(g(u))), \quad \forall u \in U.}
        \]
    \end{definition}

    \subsection{Simplessi e caratteristica di Eulero}

    \begin{definition}[$m$-simplesso]
        Un \textbf{$m$-simplesso} $\Delta^{(m)}$ in $\RR^k$ con $k \geq m$ è
        definito come l'inviluppo convesso di $m+1$ punti
        affinemente indipendenti. \smallskip

        Si dice \textbf{faccia} di $\Delta^{(m)}$ un simplesso generato
        da alcuni dei generatori di $\Delta^{(m)}$.
    \end{definition}

    \begin{definition}[Complesso simpliciale]
        Si dice \textbf{complesso simpliciale} l'unione di
        simplessi in $\RR^k$ che si intersecano a due a due
        nell'insieme vuoto oppure in una faccia.
    \end{definition}

    \begin{fact}
        Ogni varietà $M$ è omeomorfa ad un complesso simpliciale,
        finito se $M$ è compatta.
    \end{fact}

    \begin{definition}[Caratteristica di Eulero-Poincaré]
        Sia $M$ compatta. Allora si definisce la sua \textbf{caratteristica
            di Eulero-Poincaré} $\chi(M)$ come:
        \[
            \boxed{\chi(M) \defeq \sum_{i \geq 0} (-1)^i s_i(C),}
        \]
        dove $C$ è un complesso simpliciale finito a cui $M$ è omotopicamente equivalente
        e $s_i(C)$ è il numero di $i$-simplessi in $C$.
    \end{definition}

    \begin{fact}
        La caratteristica di Eulero-Poincaré è ben definita e invariante
        per equivalenza omotopica.
    \end{fact}

    \subsection{Teorema di Poincaré-Hopf}

    \begin{theorem}[Poincaré-Hopf] \label{thm:poincare_hopf}
        Sia $M$ una varietà compatta con bordo, eventualmente vuoto.
        Se $v : M \to \RR^k$ è un campo vettoriale tangente con
        zeri isolati e $\restr{v}{\partial M}$ è esterno in ogni punto
        (se $\partial M \neq \emptyset$), allora:
        \[
            \boxed{\sum_{z \in v\inv(0)} \ind(v, z) = \chi(M).}
        \]
    \end{theorem}

    \begin{corollary}
        Si può calcolare $\chi(S^m)$ nel seguente modo:
        \[
            \boxed{\chi(S^m) = \begin{cases}
                    0 & \text{se $m$ è dispari}, \\
                    1 & \text{se $m$ è pari}.
                \end{cases}}
        \]
    \end{corollary}

    \begin{proof}
        Consideriamo il campo vettoriale tangente:
        \[
            v : S^m \to \RR^{m+1}, \quad v(x) = \pi_{x^\perp}(N) = N - (N \cdot x) x,
        \]
        dove $N = e_{m+1} \in S^m$. Osserviamo che $v$ si annulla solamente in $\pm N$.
        Consideriamo le parametrizzazioni locali $g_\pm : U \to \RR^m \times \RR$ di $\pm N$, dove:
        \[
            g_\pm(u) = (u, \pm \sqrt{1 - u \cdot u}), \quad U = B_1(0) \subseteq \RR^m.
        \]
        Osserviamo che:
        \[
            \pd{}{u_i} \sqrt{1 - u \cdot u} = - \frac{u_i}{\sqrt{1 - u \cdot u}},
        \]
        da cui si deduce che:
        \[
            J g_\pm = \begin{pmatrix}
                I \\
                \mp \frac{u}{\sqrt{1 - u \cdot u}}
            \end{pmatrix}.
        \]
        Posto $\xi_\pm(u) \defeq \mp \sqrt{1 - u \cdot u} \, u$, si osserva che:
        \[
            (\dif g_\pm)_u(\xi_\pm(u)) = J (g_\pm)_u \, \xi_\pm(u) = v(g(u)).
        \]
        Quindi $\ind(v, \pm N) = \ind(\xi_\pm, 0)$. Osserviamo che $\xi_\pm$ si può
        riscalare per essere $\id_{S^{m-1}}$ nel caso positivo e la mappa
        antipodale nel caso negativo. \smallskip

        Per il Teorema \ref{thm:poincare_hopf} e il Lemma \ref{lem:grado_antipodale}, allora si ha:
        \[
            \begin{split}
                \chi(S^m) = \ind(v, N) + \ind(v, -N) = \hspace{2cm} \\
                \hspace{2cm} 1 + (-1)^m = \begin{cases}
                                              0 & \text{se $m$ è dispari}, \\
                                              2 & \text{se $m$ è pari}.
                                          \end{cases}
            \end{split}
        \]
    \end{proof}
\end{multicols*}
